\pdfinfoomitdate=1
\pdftrailerid{}
\pdfsuppressptexinfo=15
\documentclass[11pt,letterpaper]{article}
\usepackage[T1]{fontenc}
\usepackage{lmodern}
\usepackage[margin=1in]{geometry}
\usepackage{amsmath,amssymb,amsthm,mathtools}
\usepackage{booktabs,longtable,array}
\usepackage{microtype}
\usepackage[hidelinks,pdfauthor={},pdftitle={Strict partition chains and their asymptotic expansions},pdfcreator={},pdfproducer={}]{hyperref}
\usepackage{enumitem}
\setlist{itemsep=3pt,topsep=5pt}
\setlength{\emergencystretch}{2em}
\newtheorem{theorem}{Theorem}[section]
\newtheorem{lemma}[theorem]{Lemma}
\newtheorem{proposition}[theorem]{Proposition}
\newtheorem{corollary}[theorem]{Corollary}
\theoremstyle{definition}\newtheorem{definition}[theorem]{Definition}
\theoremstyle{remark}\newtheorem{remark}[theorem]{Remark}
\newcommand{\Q}{\mathbb Q}
\newcommand{\N}{\mathbb N}
\newcommand{\E}{\mathbb E}
\newcommand{\Var}{\operatorname{Var}}
\newcommand{\Log}{\operatorname{Log}}
\newcommand{\coeff}[2]{[#1^{#2}]}
\newcommand{\fall}[2]{(#1)_{#2}}
% Added by the ProveIt write phase (batch 106, 5 October 2026): dated notes
% marked [write] and a macro for file and label names.
\newcommand{\file}[1]{\nolinkurl{#1}}
\expandafter\def\expandafter\UrlBreaks\expandafter{\UrlBreaks\do\-}
\newenvironment{writenote}[1][]{\par\smallskip\begingroup\small
 \noindent\textbf{[write]}\if\relax\detokenize{#1}\relax\else~\textbf{#1.}\fi\ }%
 {\par\endgroup\smallskip}
\title{Strict partition chains and their asymptotic expansions}
\author{Report 223}
\date{}
\begin{document}
\maketitle
\begin{abstract}
The Lengyel numbers, OEIS A005121, count strict chains between the two endpoints of a partition lattice. We give a self-contained recurrence proof of their complete fixed-order asymptotic expansion on the known factorial-square scale. An elementary factorial majorant and a finite block-defect expansion yield a uniform kernel remainder. A weighted first-difference operator has limiting norm $2\log 2-1<1$, so recurrence residuals give genuine asymptotic errors and one positive, anchored leading constant. The same argument on a fixed complex marking disk proves all fixed cumulant expansions for chain length. We also give a constructive Lambert-$W$ inverse, eventual recovery on sampled values, and two-ceiling enclosures for arbitrary thresholds. Exact counting and testing code accompanies the article. The leading equivalent, its constant, the leading mean, and the central limit theorem are prior results. Formal higher corrections were previously announced; no priority claim for these coefficients or general inversion machinery is made.
\end{abstract}

\begin{writenote}[Status and trust boundary, added 5 October 2026, batch~106]
This is bundle Report~223 of an external AI-assisted research session,
filed in ProveIt as the single-source research report
\file{a005121-strict-partition-chains} (provenance, the sources as the write
read them, relation to the repository, reading conventions and the collected
non-claims in Section~\ref{spc:sec:provenance}). It is unrefereed and not
formalized; its author line reads ``Report 223'' and its PDF author field is
empty. \emph{What rests on what.} Theorems~\ref{spc:thm:main}
and~\ref{spc:thm:marked} and everything derived from them are proved by hand
in Sections~\ref{spc:sec:domination}--\ref{spc:sec:inverse} from the
recurrence~\eqref{spc:eq:recurrence}; no proof uses a computation or a
numerical value of~$C$. The decimals of Section~\ref{spc:sec:computation}
($C\approx1.098685805525187$, $K_1$, $K_2$) are uncertified diagnostics. The
leading equivalent, the positivity of~$C$, the leading mean and the central
limit theorem are re-proofs of results of Babai--Lengyel and Van
Cutsem--Ycart, as the source says. The write adds two remarks with proofs:
Remark~\ref{spc:rem:oeis} (the generating-function formula printed in OEIS
A005121 has the wrong sign; its only formal solution is the generating
function of $-z_n$) and Remark~\ref{spc:rem:transseries} (which parts of
Section~\ref{spc:sec:inverse} are instances of the transseries volume's
theorems, and a second route to Theorem~\ref{spc:thm:threshold}); and
Section~\ref{spc:sec:further}, the open questions, with a numerical
observation on the depth-sum route through
\file{a139383-iterated-bell-diagonals}.
\end{writenote}
\tableofcontents

\section{Objects and the main conclusions}
Let $\Pi_n$ be the lattice of set partitions of $[n]=\{1,\ldots,n\}$, ordered by refinement. Its least element consists of $n$ singleton blocks, and its greatest element has one block. A \emph{strict chain} here starts at the least element, ends at the greatest element, and has strictly increasing successive partitions. It need not be saturated. Its length $H$ is the number of strict transitions, not the number of partitions listed. In particular $H_1=0$; for $n\ge2$, $1\le H\le n-1$.

Let $z_n$ be the number of these chains, and let $Z_n(q)$ count them with weight $q^H$. If $S(n,j)$ denotes a Stirling number of the second kind, then
\begin{equation}\label{spc:eq:recurrence}
 Z_1(q)=1,\qquad Z_n(q)=q\sum_{j=1}^{n-1}S(n,j)Z_j(q)\quad(n\ge2),
 \qquad z_n=Z_n(1).
\end{equation}
Indeed, the first transition has a partition with $j<n$ blocks, of which there are $S(n,j)$ possibilities. Its upper interval is a partition lattice on those $j$ blocks. This proves the recurrence without any independence assumption. The first values are
\[
(z_1,\ldots,z_{10})=(1,1,4,32,436,9012,262760,10270696,518277560,32795928016).
\]
At $q=1$ choose a chain uniformly. Its probability generating function is exactly $Z_n(q)/z_n$.

Write throughout
\begin{equation}\label{spc:eq:scale}
 L=\log 2,\qquad b=1+\frac L3,\qquad
 f_n=(n!)^2(2L)^{-n}n^{-b}.
\end{equation}
\begin{theorem}[Unmarked fixed-order expansion]\label{spc:thm:main}
There is a positive constant $C$, uniquely anchored by
\begin{equation}\label{spc:eq:C}
 C=\lim_{n\to\infty}\frac{z_n}{f_n},
\end{equation}
such that for every fixed integer $M\ge0$,
\begin{equation}\label{spc:eq:expansion}
 z_n=Cf_n\left(\sum_{j=0}^{M}\frac{c_j(L)}{n^j}
                  +O_M(n^{-M-1})\right),\qquad c_0=1.
\end{equation}
The coefficients $c_j\in\Q[L]$ are uniquely and recursively computable. The first four are
\begin{align}
 c_1&=\frac L6,\label{spc:eq:cs}\\
 c_2&=-\frac{L(16L^2+45L-90)}{3240},\notag\\
 c_3&=\frac{L^2(10L^2-81L-270)}{6480},\notag\\
 c_4&=\frac{L}{146966400}\bigl(1792L^5+32400L^4+2358405L^3\notag\\[-2pt]
    &\hspace{40mm}{}-1281420L^2-8108100L-408240\bigr).\notag
\end{align}
\end{theorem}
The proof in Sections~\ref{spc:sec:domination}--\ref{spc:sec:completion} is direct and does not invoke a previously known convergence theorem. This does not make the leading limit new: it is the known Lengyel constant. The point is that the error transfer is proved, rather than inferred from formal substitution.

There is also one complex neighborhood of $q=1$ on which the same polynomials give every fixed-order marked expansion. The initial contraction disk can be taken to be $|q-1|\le1/16$; a once-only further shrink makes the analytic limiting constant nonzero. With $L(q)=\Log(q+1)-\Log q$,
\begin{equation}\label{spc:eq:markedintro}
 Z_n(q)=C(q)(n!)^2[2L(q)]^{-n}n^{-1-L(q)/3}
 \left(\sum_{j=0}^{M}\frac{c_j(L(q))}{n^j}+O_M(n^{-M-1})\right).
\end{equation}
The error inside parentheses is analytic and uniform on a fixed disk. Its fixed derivatives, on smaller disks, retain the stated order. This is a statement about a normalized remainder; differentiating the whole prefactor would introduce additional $n$-dependent factors.

In particular, for constants $K_1,K_2$ determined by derivatives of $\log C(e^t)$,
\begin{align}
 \E H_n&=\frac{n}{2L}+\frac{\log n}{6}+K_1-\frac1{12n}+O(n^{-2}),\label{spc:eq:meanintro}\\
 \Var H_n&=\frac{1-L}{4L^2}n-\frac{\log n}{12}+K_2+\frac1{24n}+O(n^{-2}).\label{spc:eq:varintro}
\end{align}
Every fixed cumulant has a full fixed-order expansion of the same kind. Section~\ref{spc:sec:inverse} gives the corresponding inverse construction and explains precisely why arbitrary threshold inversion needs an enclosure rather than the ceiling of an unqualified approximation.

\subsection{What the asymptotic assertions mean}\label{spc:sec:meaning}
All orders in this article are \emph{fixed}. For each chosen $M$ there are constants and an onset index for the indicated remainder. Neither their sizes nor the onset need be uniform as $M$ grows. We prove no convergence of the infinite formal series, Gevrey estimate, optimal truncation, exponentially small remainder, local limit theorem, Edgeworth expansion, or large-deviation principle. No interval-certified decimal value of $C$ or its marked derivatives is supplied. Exact calculations and decimal diagnostics are separated in Section~\ref{spc:sec:computation}.

\subsection{Provenance, sources and reading conventions}\label{spc:sec:provenance}
\begin{writenote}[Added 5 October 2026, batch~106]
\emph{Provenance.} This report is Report~223 of the session bundle of
Reports 1--243 that ProveIt's \file{docs/incoming} intake processed as
batch~106: the archive \file{Report223.zip} (509{,}137 bytes, SHA-256
\file{c178191805f9...cf2b8f25}, 15 files at the archive root, no wrapper
directory), arrival commit \file{60f54ea06}, placed in \file{47fc7a069}. The
text is the delivered \file{Report223.tex} (787 lines), shipped as
\file{article.tex}. The delivered 19-page PDF and the checksum manifest
\file{SHA256SUMS} are not shipped (the manifest was verified, 14 of 14
entries, at placement); the delivered README is replaced by the report
README; all three are retrievable from the arrival commit. The manuscript is
undated. Its pin is ProveIt commit \file{9f19e58de}: it inspected the
article of \file{a139383-iterated-bell-diagonals} there (blob
\file{ab5014ec}, which is still the blob at the commit of this write, so what
the source says about that report concerns its current text) and the
transseries catalogue. It carries no ``prepared for private review'' line and
no personal data. It is a single-source report, so no merge choice was made.
The write prefixed the 78 delivered labels with \file{spc:} and updated the
59 references to them; it added this subsection, the notes marked [write],
Remarks~\ref{spc:rem:oeis} and~\ref{spc:rem:transseries},
Section~\ref{spc:sec:further} and a note in one bibliography entry. No
statement, proof or number of the source was changed: the added remarks are
the last statements of their sections and the added displays are unnumbered,
so every theorem and equation keeps its delivered number.

\emph{The sources, as the write read them} (5~October 2026).
\begin{itemize}
\item OEIS A005121~\cite{OEIS}, live entry: revision~\#79 of May~30, 2026,
 the revision the source inspected, with 19 terms through
 $z_{19}=879518067472225603327860638128$. Its formula line ``E.g.f.\
 satisfies Z(z) = 1/2 * (Z(exp(z)-1) - z). [Lengyel]'' is discussed in
 Remark~\ref{spc:rem:oeis}. Its asymptotic line attributes the leading
 equivalent to Babai and Lengyel and the constant
 $C_L=1.0986858055\ldots$ (A086053) to Flajolet and Salvy.
\item \file{a139383-iterated-bell-diagonals} of this repository: the
 article's Part~II, Section~17.2 (``Other overlap and scope restrictions'';
 lines 909--910 of its \file{article.tex}) states the depth-sum identity
 $Z_n=\sum_{m\ge0}2^{-m-1}H(n,m)$ from $2Z(z)=z+Z(e^z-1)$ (the correct
 sign) and that ``a depth-summed equivalent alone does not give a pointwise
 proportional-depth equivalent''; its delivered source note
 \file{28-proportional-sources-literature.md} (lines 139--141) records the
 known Lengyel asymptotic. Read in full where cited.
\item The transseries volume
 \path{Analysis/Transseries/docs/series-and-transseries/Transseries_And_Inversion/transseries_and_inversion.tex}:
 \file{p0:def:core}, \file{p0:prop:factorial-core},
 \file{p0:thm:lambert-core}, \file{p0:thm:core-reversion},
 \file{p0:thm:backward-error}, \file{p0:def:three-inverses} and
 \file{p0:thm:staircase}, read for Remark~\ref{spc:rem:transseries}.
\item Not read by the write: Lengyel~\cite{Lengyel}, Babai and
 Lengyel~\cite{BabaiLengyel}, Prellberg~\cite{Prellberg},
 Mishna~\cite{Mishna}, Van Cutsem and Ycart~\cite{VanCutsemYcart},
 Finch~\cite{Finch} and Dickey and Rosenberg~\cite{DickeyRosenberg}. The
 statements made about them in Section~\ref{spc:sec:prior} are the
 source's; \file{data/SOURCES.json} records what the source inspected, with
 SHA-256 fingerprints of the files it read. In particular, whether
 Lengyel's 1984 paper prints the generating-function equation with the sign
 of the OEIS line was not checked.
\end{itemize}

\emph{What was checked at intake.} The batch-106 dossier read the manuscript
in full and found no error in it. Independently of the delivered code it
computed $z_n$ by the recurrence~\eqref{spc:eq:recurrence} to $n=500$:
$z_1,\dots,z_{10}$ agree with the OEIS, and
$z_n/\bigl(f_n(1+c_1/n+c_2/n^2+c_3/n^3)\bigr)$ equals
$1.0986858052222$, $1.0986858055062$, $1.0986858055214$,
$1.0986858055240$, $1.0986858055247$ at $n=100,200,300,400,500$, within
about $5\cdot10^{-13}$ of the source's $1.098685805525187$ at $n=500$, which
confirms $c_1$, $c_2$, $c_3$ and the decimal of~$C$ (an observation, not a
certificate). It checked the induction of Lemma~\ref{spc:lem:factorial} and
confirmed the OEIS sign discrepancy in the live entry. The write rechecked
identity~\eqref{spc:eq:arec} from the Stirling recurrence, the coefficient
sums $Z_n(1)=z_n$ of the six listed chain polynomials, and the OEIS history
of the formula line (Remark~\ref{spc:rem:oeis}). The delivered programs
passed on copies at placement and at the write
(Section~\ref{spc:sec:computation}, note at the end).

\emph{Relation to the repository.} Before batch~106 the repository mentioned
A005121 only in \file{a139383-iterated-bell-diagonals}, as background (the
two places above); no file mentioned A086053 or A008826. That report is not
a host for this one: it answers no question of this report, this report
answers none of its questions, and the methods differ (here a Stirling-kernel
recurrence with a first-difference contraction; there parabolic iteration of
$e^z-1$ and a fixed Fatou-coordinate contour). The source's own
Question~6 (Section~\ref{spc:sec:questions}) asks for the comparison; it is
Question~\ref{spc:q:depth} below, with a numerical observation.
\textup{[Dated note, 9~October 2026, batch~138: Part~IV of
\file{a290354-iterated-euler-diagonals} now proves the heuristic of
Question~\ref{spc:q:depth}, and answers Question~\ref{spc:q:disk} in part,
under four unrefereed inputs, two of them from
\file{a139383-iterated-bell-diagonals}; it uses this report only for the
agreement of the marked amplitude near $q=1$ and for the cumulant constants,
and nothing proved here depends on it. See the notes in those two items.]}
The source's
bibliography credits that report to ``V. Reshetnikov''; the report has no
author line (its title block reads ``A merged research report built from two
manuscripts''), so it is cited here by repository path
(note at~\cite{ProveItBell}). The inverse construction of
Section~\ref{spc:sec:inverse} and the transseries volume are compared in
Remark~\ref{spc:rem:transseries}: its Lambert core is an instance of the
volume's \file{p0:prop:factorial-core} (not of \file{p0:thm:lambert-core}).
No Lean or Rocq file of the repository treats partition-lattice chains; no
statement of this report is formalized, and its place in the collection
confers no formal status.
\end{writenote}

\begin{writenote}[What is not claimed, collected from the source]
All orders are fixed: no convergence of the formal series, no Gevrey
estimate, optimal truncation or exponentially small remainder, no local
limit theorem, Edgeworth expansion or large-deviation principle, no
distributional error estimate. No interval-certified value of $C$, $K_1$,
$K_2$ or of the marked derivatives; no numerically efficient constant $A_M$
in~\eqref{spc:eq:certificate}; the constants $A_M$, $Y_M$ of
Theorem~\ref{spc:thm:threshold} are not evaluated, and the exact inverse in
the code is bounded by its declared resources. The leading equivalent, the
constant, the leading mean and the central limit theorem are prior results;
no priority is claimed for the correction coefficients (formal higher
corrections were announced by Prellberg; the 1990 Flajolet--Salvy manuscript
was not recovered), for the inversion machinery, or for the resolution of
Van Cutsem and Ycart's general variance conjecture, which is not resolved.
The disk $|q-1|\le1/16$ is convenient, not optimal, and gives no global
marked statement. The generating function is formal (radius zero). The
source's comparison with the literature is ``an inspected source boundary,
not a worldwide-priority certificate''; that the ProveIt iterated-Bell
results cannot imply aggregate statements is not claimed. Cross-version byte
identity of the rebuild is not promised. The write adds: its numerical
observation in Question~\ref{spc:q:depth} is a floating-point extrapolation
and certifies nothing.
\end{writenote}

\medskip\noindent\begin{minipage}{\textwidth}
\emph{Reading conventions} (letters the manuscript reuses, with the tempting
false readings; no symbol was renamed).

\smallskip
\begingroup\small\renewcommand{\arraystretch}{1.1}
\noindent\begin{tabular}{@{}>{\raggedright\arraybackslash}p{0.8in}>{\raggedright\arraybackslash}p{3.4in}>{\raggedright\arraybackslash}p{1.95in}@{}}
\toprule
Symbol & Meaning here (where) & Not to be read as\\
\midrule
$z_n$, $Z_n(q)$ & strict endpoint chains of $\Pi_n$ (A005121), and their
 length polynomial; $z_n=Z_n(1)$ & a139383's $Z_n$ is the same number,
 written as a depth sum\\
$\mathcal Z(z)$ & the formal e.g.f.\ $\sum z_nz^n/n!$ (Section~\ref{spc:sec:prior}) &
 the OEIS's $Z(z)$, which obeys the misprinted equation
 (Remark~\ref{spc:rem:oeis})\\
$H$, $H_n$ & chain length (number of strict transitions; $H_1=0$) &
 a139383's iterated Bell numbers $H(n,m)$; $H(x)$ and $H_{e,k}(t)$ below\\
$H(x)$, $H_{e,k}(t)$ & $2x\log(x/e^\kappa)$ (Section~\ref{spc:sec:inverse});
 a Taylor function (Lemma~\ref{spc:lem:kernel}) & the length $H_n$\\
$L$, $L(q)$ & $\log2$; $\Log(q+1)-\Log q$, so $L(1)=L$ & a target value
 (the transseries volume writes $L$ for the target)\\
$b$ & $1+L/3$ & the coefficient $b$ of $aX+b\log X$ in
 \file{p0:thm:lambert-core}\\
$C$, $C(q)$, $C_M$, $C_e(k)$ & Lengyel's constant; its marked analytic
 version; the order-$M$ limits (all equal to $C$); block-configuration counts
 (Section~\ref{spc:sec:kernel}) & each other\\
$K(n,k)$; $K$; $K_1,K_2$ & the normalized kernel; $2\pi C$
 (Section~\ref{spc:sec:inverse}); the mean and variance constants & each
 other\\
$T_r(L)$, $T_r(n)$ & Touchard polynomials \eqref{spc:eq:touchard}; the
 weighted norm \eqref{spc:eq:T} in Lemma~\ref{spc:lem:transfer} & each other\\
$q$, $q_j(k;L)$ & the marking variable; the kernel polynomials
 \eqref{spc:eq:kernelexp} & \\
$d$, $d_n$, $d_j$ & a block count (Section~\ref{spc:sec:kernel}); first
 differences (Lemma~\ref{spc:lem:transfer}); logarithmic coefficients
 \eqref{spc:eq:d} & each other\\
$e$ & Euler's number; the block excess in Section~\ref{spc:sec:kernel} &\\
$\kappa$, $\kappa_r$ & $1+\tfrac12\log(2L)$ (Section~\ref{spc:sec:inverse});
 the $r$th cumulant & the slope $\kappa=2$ of \file{p0:prop:factorial-core}
 (Remark~\ref{spc:rem:transseries})\\
$\beta$, $\mu$ & $L/3$; the kernel mean $2L$ & a139383's slope parameter
 $\beta=n/m$\\
$A$, $B$, $R$; $A(t)$, $B(t)$, $D(t)$ & constants of Lemma~\ref{spc:lem:transfer};
 cumulant functions \eqref{spc:eq:ABD} & each other; the disks $D_0$, $D$\\
$\Lambda$, $w$, $X$, $U$, $v_m$ & $2(1+w)$; $W_0(y/(2e^\kappa))$; $e^{\kappa+w}$;
 $\log X$; inverse coefficients \eqref{spc:eq:inverseformal} & $\Lambda$ is
 the core slope of the transseries volume (same meaning)\\
$N(y)$ & $\min\{n\ge2:\log z_n\ge y\}$ & a threshold on $z_n$ itself (the
 input is $\log$-scale)\\
\bottomrule
\end{tabular}
\endgroup
\end{minipage}

\section{Prior results and the scope of this report}\label{spc:sec:prior}
Lengyel~\cite{Lengyel} proved two-sided bounds on the scale~\eqref{spc:eq:scale}; his Theorem~1.2 announced a convergence result. Babai and Lengyel~\cite{BabaiLengyel}, Theorem~2, proved the positive limit~\eqref{spc:eq:C} using a general nearly convex recurrence criterion. The end of their Section~2 explicitly asks about convergence speed. Their work also supplies the proper historical setting for the recurrence method. Near-diagonal Stirling expansions, including the fixed-deficit polynomial structure, have earlier origins; Lengyel's discussion credits Hsu.

Prellberg's 2002 presentation~\cite{Prellberg} gives the leading equivalent and a contour expression for the constant by analytic iteration. Its final slide describes additional coefficient computations as using a separate, non-rigorous method. Mishna's seminar summary~\cite{Mishna} records the leading chain asymptotic and the analytic-iteration approach. Thus a rigorous fixed-order remainder proof must be distinguished from an assertion that higher coefficients were previously unknown.

Finch~\cite{Finch}, Section~5.7, credits an unfinished 1990 Flajolet--Salvy manuscript with computation of the constant and analytic-iteration work. That full manuscript was not available for this comparison. This unresolved source prevents a claim about the first occurrence of the correction coefficients or an exhaustive priority determination.

Van Cutsem and Ycart~\cite{VanCutsemYcart}, Theorem~4.1 and Example~1, give the leading mean behavior and a central limit theorem for the strict partition-chain model. Their printed page~997 distinguishes those conclusions from a variance equivalent conjectured under their broader hypotheses. The present marked proof concerns only recurrence~\eqref{spc:eq:recurrence}. It is not a resolution of that general conjecture. Our central limit argument is a reproof of a known result; the complex-uniform fixed-order expansion is what supports the detailed cumulant corrections here.

Related proportional-depth iterated-Bell results already appear in the ProveIt article~\cite{ProveItBell}. The actual article, at the pinned revision stated in the references, was compared. It discusses A005121 as a weighted depth aggregate and proves all-order results for a different, related observable. Those results are relevant context; no claim is made that they cannot imply aggregate statements. The route here is a direct recurrence argument for the aggregate and its chain-length marking. Likewise the inverse method specializes the established Lambert-$W$ and residual-transport apparatus already represented in ProveIt's transseries work~\cite{ProveItTransseries}; it is not a new general inverse theory.

A recent phylogenetic interpretation is given by Dickey and Rosenberg~\cite{DickeyRosenberg}: the simultaneous-merger diagonal in their model is A005121. The arXiv article was compared for its enumerative scope; the published record is included bibliographically. These remarks describe an inspected source boundary, not a worldwide-priority certificate.

\subsection{A formal generating function and a checked source discrepancy}
Set $\mathcal Z(z)=\sum_{n\ge1}z_nz^n/n!$ as a \emph{formal} series. Adding the diagonal term $S(n,n)z_n=z_n$ to~\eqref{spc:eq:recurrence} gives
\begin{equation}\label{spc:eq:formal}
 2\mathcal Z(z)=z+\mathcal Z(e^z-1).
\end{equation}
The coefficient of $z$ on both sides is $2$. The OEIS raw source record~\cite{OEIS}, revision \#79 of 30 May 2026, instead displays a minus sign before $z$ in the equivalent half-formula. That sign conflicts already with $z_1=1$; equation~\eqref{spc:eq:formal} follows immediately from the recurrence. This report records the discrepancy without editing the entry.

Equation~\eqref{spc:eq:formal} is not an analytic equation for a convergent exponential generating function at the origin. The proved growth gives
\[
 \left(\frac{z_n}{n!}\right)^{1/n}\sim\frac{n}{2eL}\longrightarrow\infty,
\]
so its radius of convergence is zero. Formal composition is nevertheless well defined because $e^z-1$ has zero constant term.

\begin{remark}[{[write]} The sign in OEIS A005121, 5~October 2026]\label{spc:rem:oeis}
The live OEIS entry A005121~\cite{OEIS} (revision~\#79 of May~30, 2026, read
5~October 2026) has the formula line
\begin{center}
\texttt{E.g.f. satisfies Z(z) = 1/2 * (Z(exp(z)-1) - z). [Lengyel]}
\end{center}
The line carries no contributor's signature or date. In the entry's edit
history it first appears in revision~\#8 of September~13, 2003 (an edit by
N.~J.~A. Sloane, together with the asymptotic line), attributed
``(Lengyel)''; revision~\#30 of April~18, 2014 changed the attribution to
``[Lengyel]''; the sign has not changed since 2003. With
$\mathcal Z(z)=\sum_{n\ge1}z_nz^n/n!$ the correct equation is
\eqref{spc:eq:formal}, that is $\mathcal Z(z)=\tfrac12\bigl(\mathcal
Z(e^z-1)+z\bigr)$: the OEIS line has $-z$ where $+z$ is right.

\emph{Proof.} For any formal series $W(z)=\sum_{n\ge0}w_nz^n/n!$, the
expansion $(e^z-1)^k/k!=\sum_{n\ge k}S(n,k)z^n/n!$ gives
$W(e^z-1)=w_0+\sum_{n\ge1}\bigl(\sum_{k=1}^nS(n,k)w_k\bigr)z^n/n!$. Hence
the equation $2W(z)=W(e^z-1)+\varepsilon z$, with $\varepsilon=\pm1$, holds
if and only if $2w_0=w_0$ and
\[
 2w_n=\sum_{k=1}^{n}S(n,k)w_k+\varepsilon[n=1]\quad(n\ge1),
\]
that is
\[
 w_0=0,\qquad w_1=\varepsilon,\qquad
 w_n=\sum_{k=1}^{n-1}S(n,k)w_k\quad(n\ge2),
\]
using $S(n,n)=1$. So each sign has exactly one formal solution, and
by~\eqref{spc:eq:recurrence} at $q=1$ it is $w_n=\varepsilon z_n$. For
$\varepsilon=+1$ this is $\mathcal Z$, proving~\eqref{spc:eq:formal}. For
$\varepsilon=-1$, the OEIS line, the only solution is $-\mathcal Z$, the
generating function of $(-1,-1,-4,-32,-436,\dots)$. For $\mathcal Z$ itself
the right-hand side of the printed line is
$\tfrac12\bigl(\mathcal Z(e^z-1)-z\bigr)=\mathcal Z(z)-z$
by~\eqref{spc:eq:formal}, so the line fails exactly at the coefficient
of~$z$, where it would require $z_1=\tfrac12(z_1-1)$, that is $1=0$.
Equivalently, $-\mathcal Z$ satisfies the printed line, and $\mathcal Z$ does
not. \textup{(Wording sharpened after the independent check of 5~October
2026; the first wording read ``already the coefficient of $z$ in the OEIS
line reads $1=\tfrac12(1-1)=0$ for $\mathcal Z$''.)}
\qed

The same entry's PARI program (Michael Somos, September~22, 2007) iterates,
with series truncations, \texttt{A = x - A + subst(A, x, exp(x)-1)}, whose fixed point satisfies
$2A=x+A(e^x-1)$, the correct sign, and \file{a139383-iterated-bell-diagonals}
(Part~II, Section~17.2) prints $2Z(z)=z+Z(f(z))$ with $f(z)=e^z-1$, also
correct. The source records the discrepancy at the raw OEIS record of the
same revision (above); it is recorded here only, and nothing was submitted
to the OEIS.
\end{remark}

\section{An elementary global kernel bound}\label{spc:sec:domination}
Put $\fall nk=n(n-1)\cdots(n-k+1)$ and define
\[
 a(n,k)=\frac{2^kS(n,n-k)}{\fall nk^2}\quad(0\le k\le n-1),
 \qquad a(n,0)=1,
\]
with value zero outside the natural range when used in recurrences. The expression $\fall nk^2$ means $(\fall nk)^2$.
\begin{lemma}[Factorial domination]\label{spc:lem:factorial}
For $0\le k\le n-1$,
\begin{equation}\label{spc:eq:a-bound}
 0\le a(n,k)\le\frac1{k!}.
\end{equation}
\end{lemma}
\begin{proof}
The Stirling recurrence gives the exact identity
\begin{equation}\label{spc:eq:arec}
 a(n,k)=\frac{2(n-k)}{n^2}a(n-1,k-1)
       +\frac{(n-k)^2}{n^2}a(n-1,k).
\end{equation}
The boundary term outside the natural range is zero. Starting with $a(1,0)=1$, induction and nonnegativity imply
\[
 k!a(n,k)\le\frac{2k(n-k)+(n-k)^2}{n^2}
       =1-\frac{k^2}{n^2}\le1.
\]
The case $k=0$ is immediate.
\end{proof}
The normalized kernel and its limiting weights are
\begin{equation}\label{spc:eq:K}
 K(n,k)=S(n,n-k)\frac{f_{n-k}}{f_n}
 =L^ka(n,k)(1-k/n)^{-b},\qquad p_k=\frac{L^k}{k!}\quad(k\ge1).
\end{equation}
Lemma~\ref{spc:lem:factorial} yields
\begin{equation}\label{spc:eq:global}
 K(n,k)\le 2^bp_k\quad(k\le n/2),\qquad
 K(n,k)\le n^bp_k\quad(1\le k<n).
\end{equation}
For every fixed $\varepsilon>0$, real $A$, and integer $d\ge0$,
\begin{equation}\label{spc:eq:superpoly}
 n^A\sum_{\varepsilon\sqrt n<k<n}(1+k)^dK(n,k)\longrightarrow0.
\end{equation}
To check this explicitly, bound $(1+k)^d$ and $n^b$ by powers of $n$. For large $m$ the successive ratios in $L^k/k!$ for $k\ge m$ are at most $1/2$, so its tail is at most $2L^m/m!$. With $m\asymp\sqrt n$ and $m!\ge(m/e)^m$, this is $\exp(-c\sqrt n\log n)$ for some $c>0$, which beats every fixed inverse power. The same tail argument applies to any fixed polynomial in $k$ times $p_k$.

\section{A uniform expansion by block defects}\label{spc:sec:kernel}
A fixed-$k$ Stirling expansion alone would not justify summing a recurrence with $k$ up to $n-1$. The following construction supplies both the coefficients and the uniform remainder.

\subsection{An exact finite decomposition}
For a partition with deficit $k$, let $m_j$ be the number of its blocks of size $j\ge2$. Define
\[
 e=\sum_{j\ge3}(j-2)m_j,\qquad
 d=\sum_{j\ge3}(j-1)m_j,\qquad m_2=k-d.
\]
The number of vertices in nonsingleton blocks is $2k-e$. Counting labeled blocks gives exactly
\begin{equation}\label{spc:eq:block}
 k!a(n,k)=\sum_{e\ge0}C_e(k)\frac{\fall n{2k-e}}{(\fall nk)^2},
 \qquad
 C_e(k)=\sum_{\substack{m_j\ge0\ (j\ge3)\\\sum(j-2)m_j=e}}
 \frac{2^d\fall kd}{\prod_{j\ge3}m_j!(j!)^{m_j}}.
\end{equation}
Terms with $d>k$ vanish at integer $k$, so only legitimate block configurations contribute. For example,
\[
 C_0=1,\qquad C_1=\frac23\fall k2,\qquad
 C_2=\frac13\fall k3+\frac29\fall k4.
\]
For a nonzero term with $e\ge1$, $e+1\le d\le\min(2e,k)$. There are at most $2^e$ configurations (integer partitions of $e$ are no more numerous than compositions), and therefore
\begin{equation}\label{spc:eq:Ce}
 0\le C_e(k)\le(8k^2)^e.
\end{equation}
Moreover, for $k\le n/2$ and a legitimate numerator,
\begin{equation}\label{spc:eq:blockbound}
 \frac{\fall n{2k-e}}{(\fall nk)^2}
 \le n^{-e}\exp(2k^2/n).
\end{equation}
Indeed, use $n^{2k-e}$ for the numerator and
$-2\sum_{i<k}\log(1-i/n)\le4\sum_{i<k}i/n\le2k^2/n$
for the denominator.

\subsection{The complex Taylor estimate}
\begin{lemma}[Uniform kernel expansion]\label{spc:lem:kernel}
For every fixed integer $J\ge0$ there are polynomials
$q_j(k;L)\in\Q[L,k]$, of degree at most $2j$ in $k$, such that uniformly for $1\le k\le\sqrt n/8$,
\begin{equation}\label{spc:eq:kernelexp}
 K(n,k)=p_k\left(\sum_{j=0}^Jq_j(k;L)n^{-j}
       +O_J\bigl((1+k)^{2J+2}n^{-J-1}\bigr)\right).
\end{equation}
Here $q_0=1$ and $q_j(0;L)=0$ for $j\ge1$. The same result is uniform when $L$ ranges over a fixed compact complex neighborhood of $\log2$ and $b=1+L/3$.
\end{lemma}
\begin{proof}
For every relevant $e$, set
\[
 H_{e,k}(t)=\frac{\prod_{i=0}^{2k-e-1}(1-it)}
                 {\prod_{i=0}^{k-1}(1-it)^2}(1-kt)^{-b},
 \qquad R_k=\frac1{8k^2}.
\]
The branches at zero are used. On $|t|\le R_k$ all factors are nonzero, and their arguments $it$ have modulus at most $1/4$. Since
$|\log(1-z)|\le(4/3)|z|$ for $|z|\le1/4$, summing the logarithms gives
\[
 |\log H_{e,k}(t)|\le\frac{3+|b|}{6},\qquad
 |H_{e,k}(t)|\le B:=\exp((3+|b|)/6).
\]
For clarity, the sum of numerator indices is at most $2k^2$, twice the denominator sum is at most $k^2$, and the marking term contributes at most $|b|k$. This proves the stated bound uniformly in the relevant $e$ and $k$.

Write $H_{e,k}(t)=\sum_{r\ge0}A_{e,r}(k)t^r$. Since $8k^2/n\le1/8$, Cauchy's estimate yields, for every fixed $D\ge0$,
\begin{equation}\label{spc:eq:cauchy}
 \left|H_{e,k}(1/n)-\sum_{r=0}^DA_{e,r}(k)n^{-r}\right|
 \le\frac{8B}{7}\,8^{D+1}k^{2D+2}n^{-D-1}.
\end{equation}
For $e\le J$ take $D=J-e$ and multiply by $C_e(k)n^{-e}$. Equation~\eqref{spc:eq:Ce} makes every resulting remainder
$O_J(k^{2J+2}n^{-J-1})$. The remaining defect tail is bounded by
\[
 \sum_{e\ge J+1}C_e(k)\frac{\fall n{2k-e}}{(\fall nk)^2}
 \le e^{2/64}\frac{(8k^2/n)^{J+1}}{1-8k^2/n}.
\]
The extra factor $(1-k/n)^{-b}$ is uniformly bounded, completing the remainder proof.

Finally, for $r\ge1$ the coefficient of $t^r$ in $\log H_{e,k}$ is
\begin{equation}\label{spc:eq:logH}
 g_{e,r}(k)=\frac{2\sum_{i=0}^{k-1}i^r-
                    \sum_{i=0}^{2k-e-1}i^r+bk^r}{r}.
\end{equation}
For fixed $e$ this is a polynomial in $k$ of degree at most $r+1$. Exponentiation gives $\deg_k A_{e,r}\le2r$, while $\deg C_e\le2e$. Thus
\begin{equation}\label{spc:eq:qconstruct}
 q_j(k;L)=\sum_{e=0}^j C_e(k)A_{e,j-e}(k)
\end{equation}
has degree at most $2j$. Polynomial power sums may be used when defining these polynomials: terms that would have a negative falling-factorial length at a small integer $k$ have $C_e(k)=0$ and contribute nothing. The finitely many such small $k$ can equivalently be handled directly. For $e>0$, $C_e$ has a factor $k$; the $e=0,k=0$ series is $1$. Hence $q_j(0;L)=0$ for $j>0$. All bounds use only a bound for $|b|$, proving the complex-uniform version as well.
\end{proof}
The first coefficients are
\begin{equation}\label{spc:eq:q12}
 \begin{split}
 q_1(k;L)&=\frac{k(L-k+1)}3,\\
 q_2(k;L)&=\frac{k^2}{18}\bigl(L^2-2Lk+5L+k^2-6k+5\bigr).
 \end{split}
\end{equation}
Equations~\eqref{spc:eq:global} and~\eqref{spc:eq:superpoly} allow the expansion to be summed term by term through any fixed order with any fixed polynomial weight in $k$. In particular,
\begin{equation}\label{spc:eq:moments}
 \sum_{k\ge1}p_k=1,\quad \mu:=\sum_{k\ge1}kp_k=2L,
 \quad\sum_{k\ge1}p_kq_1(k;L)=0,\quad
 \sum_{k\ge1}p_kq_2(k;L)=-\frac{L^2}3.
\end{equation}

\section{A triangular construction of every coefficient}\label{spc:sec:coefficients}
Set $P_M(n)=\sum_{s=0}^Mc_sn^{-s}$ with $c_0=1$. We choose the coefficients so that
\begin{equation}\label{spc:eq:residual}
 \sum_{k=1}^{n-1}K(n,k)P_M(n-k)-P_M(n)=O_M(n^{-M-2}).
\end{equation}
To see why this is constructive, for $s\ge1$ expand
\[
 (n-k)^{-s}=n^{-s}\sum_{h\ge0}\binom{s+h-1}{h}(k/n)^h
\]
only through the required finite order; for $s=0$ just the $h=0$ term is present. On the near range this finite Taylor expansion has a uniform remainder; on the far range~\eqref{spc:eq:global} controls the bounded finite sequence values after the patch described below. Lemma~\ref{spc:lem:kernel} through order $M+1$ then gives a residual expansion through order $n^{-M-1}$ with error $O_M(n^{-M-2})$.

The constant term vanishes because $\sum p_k=1$. The coefficient of $n^{-1}$ vanishes because the $q_1$ mean is zero. At order $n^{-m-1}$ the first uncancelled occurrence of $c_m$ has multiplier
\begin{equation}\label{spc:eq:triangular}
 \sum_{k\ge1}p_k\bigl(q_1(k;L)+mk\bigr)=m\mu=2mL\ne0.
\end{equation}
Thus $c_m$ is minus the contribution of the already known coefficients, divided by $2mL$. The contribution of $c_{m+1}$ in a longer formal series would cancel between the leading kernel and the subtracted $P_M(n)$. There is consequently no hidden higher unknown in this step. At the first nontrivial order the residual is $-L^2/3+2Lc_1$, proving $c_1=L/6$.

For an explicit finite algebraic recipe, define the moment functional on polynomials in $k$ by
\begin{equation}\label{spc:eq:touchard}
 \mathcal M(1)=1,\qquad
 \mathcal M(k^r)=2T_r(L)\quad(r\ge1),\qquad
 T_r(L)=\sum_{h=0}^rS(r,h)L^h.
\end{equation}
Form $q_0,\ldots,q_{m+1}$ using~\eqref{spc:eq:logH}--\eqref{spc:eq:qconstruct}; finite exponential coefficients can be generated by
\[
 A_{e,0}=1,\qquad
 A_{e,r}=\frac1r\sum_{a=1}^r a\,g_{e,a}A_{e,r-a}.
\]
Then take the coefficient of $t^{m+1}$ in
\begin{equation}\label{spc:eq:residualalgorithm}
 \mathcal M\left[\left(\sum_{a=0}^{m+1}q_a(k;L)t^a\right)
       \left(\sum_{s=0}^{m}c_st^s(1-kt)^{-s}\right)\right]
       -\sum_{s=0}^{m}c_st^s,
\end{equation}
and set it to zero. Every operation is finite and rational.

\begin{proposition}\label{spc:prop:ring}
The construction gives $c_m\in\Q[L]$ at every order, not merely $\Q(L)$.
\end{proposition}
\begin{proof}
Each nontrivial polynomial contribution to a positive-order residual vanishes at $k=0$: it contains either $q_a$ for $a>0$ or a positive shift power $k^h$. Every $T_r(L)$ for $r>0$ is divisible by $L$. Hence the known residual contribution is divisible by $L$ whenever earlier $c_s$ lie in $\Q[L]$. Division by $2mL$ in~\eqref{spc:eq:triangular} stays in $\Q[L]$. Induct from $c_0=1$.
\end{proof}
This produces~\eqref{spc:eq:cs} and every further fixed coefficient. In subsequent recurrence arguments, redefine $P_M(n)=1$ for the finitely many small integer indices where necessary. Since $P_M(n)\to1$, this makes it positive and bounded above and below. A change at fixed $n-k$ has $k\ge n-O(1)$, so~\eqref{spc:eq:global} makes its effect smaller than every inverse power. The residual~\eqref{spc:eq:residual} is unchanged asymptotically.

\section{Strict contraction of first differences}\label{spc:sec:contraction}
The formal coefficient calculation is not yet an asymptotic theorem. The following transfer lemma supplies the needed stability and limiting constant.
\begin{lemma}[Positive first-difference transfer]\label{spc:lem:transfer}
Suppose $u_n>0$ satisfies
\[
 u_n=\sum_{k=1}^{n-1}w(n,k)u_{n-k},\qquad w(n,k)\ge0,
 \qquad \sigma_n:=\sum_kw(n,k)=1+O(n^{-r}),\quad r>1.
\]
Suppose, for fixed constants $A,B,R>0$,
\[
 w(n,k)\le A\frac{R^k}{k!}\ (k\le n/2),\qquad
 w(n,k)\le An^B\frac{R^k}{k!}\ (k<n),
\]
and $w(n,k)\to p_k$ at every fixed $k$, where $\sum p_k=1$ and $\mu=\sum kp_k<2$. Then $u_n$ is bounded and bounded away from zero, converges to some $C>0$, and
\[
 u_n=C+O(n^{1-r}).
\]
\end{lemma}
\begin{proof}
Let $U_n=\max_{j\le n}u_j$ and $m_n=\min_{j\le n}u_j$. Nonnegativity gives
\[
 U_n\le U_{n-1}\max(1,\sigma_n),\qquad
 m_n\ge m_{n-1}\min(1,\sigma_n).
\]
All row sums are positive. Beyond a finite index $|\sigma_n-1|\le1/2$, and $\sum|\sigma_n-1|<\infty$. Thus the upper product is finite and the lower product strictly positive. We have $0<a\le u_n\le A_0<\infty$ without assuming convergence.

For $n\ge2$ set $d_n=u_n-u_{n-1}$ and, for $1\le j\le n-2$,
$W(n,j)=\sum_{k=j+1}^{n-1}w(n,k)$. Telescoping the individual term $u_{n-k}$ gives the exact identity
\begin{equation}\label{spc:eq:diff}
 d_n=(\sigma_n-1)u_{n-1}-\sum_{j=1}^{n-2}W(n,j)d_{n-j}.
\end{equation}
All difference indices are at least $2$. Interchanging finite sums, its weighted norm is
\begin{equation}\label{spc:eq:T}
 T_r(n)=\sum_{j=1}^{n-2}W(n,j)\left(\frac n{n-j}\right)^r
 =\sum_{k=1}^{n-1}w(n,k)\sum_{j=1}^{k-1}\left(\frac n{n-j}\right)^r.
\end{equation}
For $k\le n/2$ the inner sum is at most $2^rk$, giving a summable majorant $\mathrm{const}\,kR^k/k!$. For $k>n/2$ it is at most $n^{r+1}$, and the factorial tail makes that entire contribution tend to zero faster than any fixed inverse power. Fixed-$k$ convergence therefore proves
\begin{equation}\label{spc:eq:Tlimit}
 T_r(n)\longrightarrow\sum_{k\ge1}(k-1)p_k=\mu-1<1.
\end{equation}
Choose $\theta<1$ and $N$ with $T_r(n)\le\theta$ for $n\ge N$. Boundedness of $u$ bounds the forcing term in~\eqref{spc:eq:diff} by $B_0n^{-r}$. Choose $A_1$ no smaller than $B_0/(1-\theta)$ and all finitely many initial $n^r|d_n|$. Induction in~\eqref{spc:eq:diff} gives
\[
 |d_n|\le n^{-r}\bigl(B_0+A_1T_r(n)\bigr)\le A_1n^{-r}.
\]
The differences are absolutely summable, so a limit $C$ exists, and
\begin{equation}\label{spc:eq:Ctail}
 |u_N-C|\le A_1\sum_{n>N}n^{-r}\le\frac{A_1}{r-1}N^{1-r}.
\end{equation}
The lower bound on $u_n$ makes $C>0$.
\end{proof}
For the Lengyel kernel, $\mu=2L$ and the limiting norm is
\[
 2\log2-1=0.3862943611\ldots<1.
\]
The inequality, not the displayed decimal, is what the proof needs. Every weighted limit concerns a fixed $r$; no uniformity for $r=r(n)$ is being used.

\section{Completion and anchoring of the expansion}\label{spc:sec:completion}
\begin{proof}[Proof of Theorem~\ref{spc:thm:main}]
For a fixed $M$, use the positively patched $P_M$ and put
\[
 h_n=f_nP_M(n),\quad u_n=\frac{z_n}{h_n},\quad
 w_M(n,k)=K(n,k)\frac{P_M(n-k)}{P_M(n)}.
\]
These weights are nonnegative and inherit the factorial bounds up to constants. They converge to $p_k$ for each fixed $k$. Dividing~\eqref{spc:eq:residual} by $P_M(n)$ gives
$\sigma_n=1+O_M(n^{-M-2})$.
Lemma~\ref{spc:lem:transfer} with $r=M+2$ yields
\[
 \frac{z_n}{f_nP_M(n)}=C_M+O_M(n^{-M-1}),\qquad C_M>0.
\]
Since $P_M(n)\to1$, each $C_M$ is the same limit of $z_n/f_n$. In particular they all equal the $M=0$ constant $C$. Multiplication by the bounded $P_M$ proves~\eqref{spc:eq:expansion}. The normalization $c_0=1$ and nonzero triangular multipliers prove uniqueness of all coefficients; the limiting normalization proves uniqueness of $C$.
\end{proof}
This also explains why successive orders do not fit independent leading constants. The constant is the same already known Lengyel limit at every order.

The proof is effective in principle. If explicit bounds and a finite computation establish $|d_n|\le A_Mn^{-M-2}$ beyond an index, then
\begin{equation}\label{spc:eq:certificate}
 \left|C-\frac{z_N}{f_NP_M(N)}\right|
 \le\frac{A_M}{M+1}N^{-M-1}.
\end{equation}
The kernel remainder gives computable bounds, and finite-cutoff estimates make the weighted contraction effective. No numerically efficient $A_M$ or evaluated interval certificate is claimed here. In particular a stable decimal extrapolation is not a substitute for~\eqref{spc:eq:certificate} with verified inputs.

\section{One complex marking disk for all fixed orders}\label{spc:sec:marked}
We now prove~\eqref{spc:eq:markedintro}. Positivity no longer holds for complex $q$; a separate boundedness bootstrap is necessary.
\begin{theorem}[Complex-uniform expansion]\label{spc:thm:marked}
On a neighborhood of the closed disk $D_0=\{q:|q-1|\le1/16\}$ define
\[
 L(q)=\Log(q+1)-\Log q,\quad b(q)=1+L(q)/3,\quad
 f_n(q)=(n!)^2[2L(q)]^{-n}n^{-b(q)},
\]
using analytic logarithms in the right half-plane. There is an analytic limit
$C(q)=\lim_n Z_n(q)/f_n(q)$ in the interior of $D_0$, with uniform convergence on $D_0$. On one smaller closed disk $D$ about $1$, $C$ is nonzero and, for every fixed $M\ge0$,
\begin{equation}\label{spc:eq:marked}
 \frac{Z_n(q)}{C(q)f_n(q)}
   =\sum_{j=0}^M c_j(L(q))n^{-j}+O_{D,M}(n^{-M-1})
\end{equation}
uniformly. The same disk $D$ works for every fixed $M$. Constants, finite patches, and onset indices may depend on $M$. On every smaller disk, each fixed derivative of the analytic remainder in~\eqref{spc:eq:marked} has the same order.
\end{theorem}
\begin{proof}
First, the explicit disk has an absolute contraction margin. Since
$L'(q)=-1/[q(q+1)]$, integration along a radial segment gives
\[
 |L(q)-\log2|\le\frac{16}{465}\quad(q\in D_0).
\]
The elementary inequality $\log2<7/10$ implies $|L(q)|<3/4$. (Already the first five terms of the exponential series show $e^{7/10}>2$.) The limiting weights are
\begin{equation}\label{spc:eq:pkq}
 p_k(q)=q\frac{L(q)^k}{k!},\quad
 \sum_{k\ge1}p_k(q)=1,\quad
 \sum_{k\ge1}kp_k(q)=(q+1)L(q).
\end{equation}
For $g(x)=(x-1)e^x+1$, $g'(x)=xe^x\ge0$, and $e^{3/4}>2$. Therefore
\begin{equation}\label{spc:eq:diskcontract}
 \sum_{k\ge1}(k-1)|p_k(q)|
 =|q|g(|L(q)|)<\frac{17}{16}g(3/4)<\frac{17}{32}<1.
\end{equation}
This disk selection contains no reference to $M$.

The kernel polynomials of Lemma~\ref{spc:lem:kernel} depend only on $L$ and not separately on $q$. For $r\ge1$ their moments now use
\begin{equation}\label{spc:eq:markedmoments}
 \sum_{k\ge1}p_k(q)k^r=(q+1)T_r(L(q)),
\end{equation}
whereas the zeroth moment remains $1$. The $q_1$ moment is identically zero. Every nontrivial residual polynomial vanishes at $k=0$; its moment therefore has a factor $(q+1)L(q)$. The triangular multiplier is $m(q+1)L(q)$. These common factors cancel, leaving exactly the same $c_m\in\Q[L]$ from Section~\ref{spc:sec:coefficients}. This proves universality at every order, not just for the displayed coefficients.

For fixed $M$, the polynomial
$P_M(n,q)=\sum_{j=0}^Mc_j(L(q))n^{-j}$ tends uniformly to $1$ on $D_0$. Choose an $M$-dependent $N_M$ beyond which $|P_M-1|\le1/2$, and set $P_M=1$ at the finitely many earlier integer indices. The patch is analytic in $q$ and bounded above and away from zero. Put
\[
 h_n(q)=f_n(q)P_M(n,q),\quad u_n(q)=Z_n(q)/h_n(q),\quad
 w(n,k;q)=qS(n,n-k)\frac{h_{n-k}(q)}{h_n(q)}.
\]
Lemma~\ref{spc:lem:factorial} and uniform bounds for the patch give constants $A,B,R$ such that
\begin{equation}\label{spc:eq:complexbounds}
 |w(n,k;q)|\le A\frac{R^k}{k!}\ (k\le n/2),\qquad
 |w(n,k;q)|\le An^B\frac{R^k}{k!}\ (k<n),
\end{equation}
uniformly on $D_0$. One may take $R=\sup_{D_0}|L(q)|$ and absorb $|q|$ into $A$. Fixed-$k$ convergence to $p_k(q)$ is uniform. The uniform Cauchy proof in Lemma~\ref{spc:lem:kernel}, the moment cancellations, and the factorial tails give
\begin{equation}\label{spc:eq:complexrow}
 \sigma_n(q):=\sum_kw(n,k;q)=1+O_{M,D_0}(n^{-r}),\qquad r=M+2.
\end{equation}
The finite patch changes only factorially small tails.

Set $W(n,j;q)=\sum_{k>j}w(n,k;q)$. The triangle inequality, followed by exactly the near/far split in the proof of~\eqref{spc:eq:Tlimit}, yields
\begin{equation}\label{spc:eq:complexnorm}
 \limsup_{n\to\infty}\sup_{q\in D_0}
 \sum_{j=1}^{n-2}|W(n,j;q)|\left(\frac n{n-j}\right)^r
 \le\sup_{q\in D_0}\sum_{k\ge1}(k-1)|p_k(q)|<1.
\end{equation}
A finite-$k$ cutoff makes the uniform convergence explicit; the remaining terms are dominated by a constant multiple of $kR^k/k!$, with the long-jump portion factorially negligible. Thus for this fixed $r$ there are $N$ and $\theta<1$ bounding every weighted norm by $\theta$ beyond $N$.

It remains to prove boundedness without positivity. Identity~\eqref{spc:eq:diff} remains valid for the complex functions. Let $B_0$ bound $n^r|\sigma_n-1|$ uniformly beyond $N$, and define
\[
 U_n=\max_{1\le j\le n}\sup_{D_0}|u_j|,\qquad
 D_n=\max_{2\le j\le n}j^r\sup_{D_0}|u_j-u_{j-1}|.
\]
If the maximum defining $D_n$ occurs after the initial segment,~\eqref{spc:eq:diff} implies $D_n\le B_0U_n+\theta D_n$. If it occurs earlier it is $D_{N-1}$. Hence in all cases
\begin{equation}\label{spc:eq:bootstrap1}
 D_n\le D_{N-1}+\frac{B_0}{1-\theta}U_n.
\end{equation}
Independently, telescoping gives, with $S_N=\sum_{j\ge N}j^{-r}$,
\begin{equation}\label{spc:eq:bootstrap2}
 U_n\le U_{N-1}+S_ND_n.
\end{equation}
Increase $N$ while retaining the already obtained $B_0,\theta$ until $B_0S_N/(1-\theta)\le1/2$. Combining the two inequalities proves
\[
 U_n\le2(U_{N-1}+S_ND_{N-1})
\]
for every $n$, and then also bounds $D_n$. The initial maxima are finite because they involve finitely many analytic functions on a compact disk. No boundedness of $u$ was assumed to obtain this conclusion.

We have proved uniformly that $|u_n-u_{n-1}|\le A_Mn^{-M-2}$, so
\[
 u_n(q)=C_M(q)+O_{M,D_0}(n^{-M-1}).
\]
Uniform convergence makes $C_M$ holomorphic in the disk interior. Since $P_M\to1$ uniformly, all $C_M$ equal the same $M=0$ limit $C$. Theorem~\ref{spc:thm:main} gives $C(1)>0$. Continuity supplies a once-only smaller disk $D$ on which $C$ has no zeros. Dividing by this bounded-away-from-zero function proves~\eqref{spc:eq:marked}. Cauchy's estimates on nested disks give the derivative assertion.
\end{proof}
The factor $n^{-b(q)}$ always means $\exp(-b(q)\log n)$ with the real logarithm of the positive integer $n$. The factor $[2L(q)]^{-n}$ has an integer exponent, so it creates no additional branch choice. The order-independent disk does not mean order-independent estimates: the eventual contraction index and the patch can grow with $M$.

\section{Chain length and all fixed cumulants}\label{spc:sec:cumulants}
Let $\ell_j(L)$ be the formal logarithmic coefficients determined by
\begin{equation}\label{spc:eq:ell}
 \log\left(\sum_{j\ge0}c_j(L)x^j\right)=\sum_{j\ge1}\ell_j(L)x^j.
\end{equation}
Only finitely many coefficients are used at any one order. In particular,
\[
 \ell_1(L)=\frac L6,\qquad
 \ell_2(L)=-\frac{L(8L^2+45L-45)}{1620}.
\]
On a sufficiently small fixed complex $t$-disk define
\begin{equation}\label{spc:eq:ABD}
 A(t)=\log\frac{L}{L(e^t)},\qquad
 B(t)=\frac{L-L(e^t)}3,\qquad
 D(t)=\log\frac{C(e^t)}{C(1)},
\end{equation}
with logarithms zero at $t=0$. Theorem~\ref{spc:thm:marked} makes the normalized bracket uniformly close to $1$ for sufficiently large $n$; its analytic logarithm therefore has the usual fixed-order Taylor remainder. If a common zero-free disk for every finite $n$ is desired, shrink once more for the finitely many remaining polynomials, which are nonzero at $q=1$.

\begin{corollary}[Logarithmic generating function]\label{spc:cor:logmgf}
For every fixed $M\ge0$, uniformly on a fixed small complex disk,
\begin{align}\label{spc:eq:logmgf}
 \log\E e^{tH_n}
 &=nA(t)+(\log n)B(t)+D(t)\notag\\
 &\quad+\sum_{j=1}^M\bigl(\ell_j(L(e^t))-\ell_j(L)\bigr)n^{-j}
       +O_M(n^{-M-1}).
\end{align}
The remainder is analytic; every fixed derivative has the same order on smaller disks. Thus the $r$th cumulant, for each fixed $r\ge1$, is
\begin{equation}\label{spc:eq:cumulants}
 \kappa_r(H_n)=nA^{(r)}(0)+(\log n)B^{(r)}(0)+D^{(r)}(0)
 +\sum_{j=1}^M\left.\frac{d^r}{dt^r}\ell_j(L(e^t))\right|_{t=0}n^{-j}
 +O_{M,r}(n^{-M-1}).
\end{equation}
\end{corollary}
\begin{proof}
Divide the marked expansion by its value at $q=1$ and take the analytic logarithm. The factorials cancel, the exponential scale gives $nA$, and the power of $n$ gives $(\log n)B$. Finite logarithmic Taylor expansion supplies the $\ell_j$ terms with a uniform analytic remainder. Apply Cauchy's formula on a nested disk to differentiate that remainder.
\end{proof}
Direct differentiation gives
\[
 A'(0)=\frac1{2L},\quad A''(0)=\frac{1-L}{4L^2},\quad
 B'(0)=\frac16,\quad B''(0)=-\frac1{12},
\]
\[
 \left.\frac d{dt}\ell_1(L(e^t))\right|_0=-\frac1{12},\qquad
 \left.\frac {d^2}{dt^2}\ell_1(L(e^t))\right|_0=\frac1{24}.
\]
Taking $K_1=D'(0)$ and $K_2=D''(0)$ proves~\eqref{spc:eq:meanintro}--\eqref{spc:eq:varintro}. These constants depend on the strict-transition convention $H_1=0$. Counting partitions in the chain instead would add $1$ to the mean and its constant.

\begin{corollary}[The known central limit theorem]\label{spc:cor:clt}
Put $v=(1-L)/(4L^2)>0$. Then
\[
 \frac{H_n-\E H_n}{\sqrt{nv}}\ \Longrightarrow\ \mathcal N(0,1).
\]
The denominator may also be replaced by the true standard deviation.
\end{corollary}
\begin{proof}
In~\eqref{spc:eq:logmgf} take $t=is/\sqrt{nv}$ for fixed real $s$ and subtract the exact mean term. The quadratic contribution of $nA$ tends to $-s^2/2$; its cubic and higher terms are $O(n|t|^3)=O(n^{-1/2})$. After centering, the analytic terms multiplied by $\log n$ or bounded coefficients contribute $o(1)$. Characteristic functions converge to $e^{-s^2/2}$, proving the limit. Equation~\eqref{spc:eq:varintro} gives $\Var H_n/(nv)\to1$.
\end{proof}
This deduction credits the earlier theorem of Van Cutsem--Ycart~\cite{VanCutsemYcart}. It is not accompanied here by a distributional error estimate or a local limit assertion.

\section{Logarithmic asymptotics and inversion}\label{spc:sec:inverse}
Define
\[
 \beta=L/3,\qquad K=2\pi C,\qquad
 \kappa=1+\tfrac12\log(2L),\qquad H(x)=2x\log(x/e^\kappa).
\]
\begin{proposition}[The logarithmic coefficients]\label{spc:prop:log}
For every fixed $M\ge0$,
\begin{equation}\label{spc:eq:logzn}
 \log z_n=H(n)-\beta\log n+\log K+
          \sum_{j=1}^Md_jn^{-j}+O_M(n^{-M-1}).
\end{equation}
The $d_j$ are obtained by adding the formal series in~\eqref{spc:eq:ell} to the odd Stirling series:
\begin{equation}\label{spc:eq:d}
 \sum_{j\ge1}d_jx^j
 =\sum_{j\ge1}\ell_j(L)x^j+
   \sum_{r\ge1}\frac{B_{2r}}{r(2r-1)}x^{2r-1}.
\end{equation}
Here the Bernoulli numbers use $B_2=1/6$. In particular,
\[
 d_1=\frac{1+L}{6},\qquad
 d_2=-\frac{L(16L^2+90L-90)}{3240}.
\]
\end{proposition}
\begin{proof}
At each fixed order, take the logarithm of the positive normalized expansion~\eqref{spc:eq:expansion}; Taylor's remainder is legitimate because the bracket tends to $1$. Add the fixed-order Stirling expansion
\[
 2\log(n!)=2n\log n-2n+\log n+\log(2\pi)
 +\sum_{r\ge1}\frac{B_{2r}}{r(2r-1)}n^{-(2r-1)},
\]
where the last series is truncated as needed with its fixed-order remainder. The $+\log n$ cancels the $-\log n$ in $-b\log n$, leaving exactly $-\beta\log n$. The remaining terms give~\eqref{spc:eq:logzn}--\eqref{spc:eq:d}.
\end{proof}

\subsection{A Lambert core and a triangular inverse}
For $y\to+\infty$ let
\begin{equation}\label{spc:eq:core}
 w=W_0\!\left(\frac{y}{2e^\kappa}\right),\qquad
 X=\frac{y}{2w},\qquad U=\log X,\qquad
 \Lambda=2(1+w)=2U-\log(2L).
\end{equation}
The principal real branch is used. Since $we^w=y/(2e^\kappa)$, we have $X=e^{\kappa+w}$, $H(X)=y$, and $H'(X)=\Lambda>0$.

For every fixed $M$ consider the smooth model
\begin{equation}\label{spc:eq:FM}
 F_M(x)=H(x)-\beta\log x+\log K+\sum_{j=1}^Md_jx^{-j}.
\end{equation}
It is strictly increasing for all sufficiently large $x$, since $F_M'(x)\sim2\log x$. Let $\xi_M(y)$ be its unique sufficiently large solution of $F_M(x)=y$.

\begin{theorem}[Constructive fixed-order inverse]\label{spc:thm:inverse}
Set $v_0(U)=(\beta U-\log K)/\Lambda$. The coefficients $v_m(U)$ for $m\ge1$ are uniquely determined by requiring the following formal expression, with $t=1/X$ and $\delta=\sum_{m\ge0}v_m(U)t^m$, to vanish:
\begin{align}\label{spc:eq:inverseformal}
 0={}&\Lambda\delta-\beta U+\log K
 +\sum_{r\ge2}\frac{2(-1)^r}{r(r-1)}\delta^rt^{r-1}\notag\\
 &-\beta\log(1+t\delta)
 +\sum_{j\ge1}d_jt^j(1+t\delta)^{-j}.
\end{align}
At order $m\ge1$, take minus the coefficient of $t^m$ in all terms except $\Lambda\delta$, using $v_0,\ldots,v_{m-1}$, and divide by $\Lambda$. Then
\begin{equation}\label{spc:eq:inverseexp}
 \xi_M(y)=X+\sum_{m=0}^Mv_m(\log X)X^{-m}
                 +O_M\!\left(\frac{X^{-M-1}}{\log X}\right).
\end{equation}
Each $v_m$ is a rational function of $U$ over
$\Q[L,\log K,\log(2L)]$; $v_0=O(1)$ and $v_m=O(1/U)$ for $m\ge1$.
\end{theorem}
\begin{proof}
Taylor expansion at $X$ gives
\[
 \frac{H^{(r)}(X)}{r!}=\frac{2(-1)^r}{r(r-1)X^{r-1}}\quad(r\ge2),
\]
which produces the first nonlinear sum in~\eqref{spc:eq:inverseformal}. The other terms follow from the Taylor expansions of $\log(X+\delta)$ and $(X+\delta)^{-j}$. Every nonlinear term carries at least one factor of $t$; hence $v_m$ first appears as $\Lambda v_m$ at degree $m$. This proves triangularity, uniqueness, and the rational-function claim. Since $\Lambda\asymp U$, $v_0$ is bounded. Inductively the remaining degree-$m$ expression is bounded for large $U$, proving $v_m=O(1/U)$.

For fixed $M$, substitute the finite displacement in $F_M$. The recursion cancels the residual through $X^{-M}$, and the remaining Taylor remainder is $O_M(X^{-M-1})$. All finitely many displacement coefficients are bounded; for large $X$ the Taylor expansions are uniform on $|\delta|/X<1/2$. On a fixed-size neighborhood of $X$, $F_M'(x)\sim2\log X$. The intermediate value theorem supplies a root there, eventual monotonicity makes it the unique large root, and the mean value theorem divides the residual by this derivative, proving~\eqref{spc:eq:inverseexp}.
\end{proof}
The first terms are
\begin{align}\label{spc:eq:v12}
 v_1&=-\frac{v_0^2-\beta v_0+d_1}{\Lambda},\notag\\
 v_2&=-\frac{(2v_0-\beta)v_1-v_0^3/3+\beta v_0^2/2-d_1v_0+d_2}{\Lambda}.
\end{align}
In particular
\[
 \xi_1(y)=X+v_0-\frac{v_0^2-\beta v_0+d_1}{X\Lambda}
              +O\!\left(\frac1{X^2\log X}\right).
\]
Although $v_0\to\beta/2=L/6$, replacing the full $U$-dependent $v_0$ by that limit discards a generally much larger $1/\log X$ correction and is not valid at the displayed precision.

\subsection{Sampled values and the threshold staircase}
The sequence is strictly increasing from $n=2$ onward, because
\begin{equation}\label{spc:eq:increasing}
 z_{n+1}\ge S(n+1,n)z_n=\binom{n+1}{2}z_n>z_n\quad(n\ge2).
\end{equation}
The equality $z_1=z_2=1$ is why range inversion is stated for $n\ge2$.
\begin{corollary}[Eventual recovery on sampled inputs]\label{spc:cor:sampled}
If the input is known to be $y=\log z_n$, then
\[
 \xi_M(y)=n+O_M\!\left(\frac{n^{-M-1}}{\log n}\right).
\]
The finite expression in~\eqref{spc:eq:inverseexp} has the same accuracy. Nearest-integer rounding therefore recovers $n$ for all sufficiently large sampled inputs.
\end{corollary}
\begin{proof}
Equation~\eqref{spc:eq:logzn} gives $F_M(n)=y+O_M(n^{-M-1})$. Around $n$ the derivative is comparable to $\log n$. Apply the same mean-value residual argument as above; also $X\sim n$. The error eventually falls below $1/2$.
\end{proof}
This is an asymptotic recovery theorem, not an implemented certified cutoff using approximate $C$.

For an arbitrary real threshold define
\[
 N(y)=\min\{n\ge2:\log z_n\ge y\}.
\]
\begin{theorem}[Two-ceiling threshold enclosure]\label{spc:thm:threshold}
For every fixed $M$ there are $A_M>0$ and $Y_M$ such that for $y\ge Y_M$, writing $\xi=\xi_M(y)$ and
\[
 \varepsilon_M(y)=A_M\frac{\xi^{-M-1}}{\log\xi},
\]
one has
\begin{equation}\label{spc:eq:ceilings}
 \left\lceil\xi-\varepsilon_M(y)\right\rceil
 \le N(y)\le
 \left\lceil\xi+\varepsilon_M(y)\right\rceil.
\end{equation}
The center can be replaced by the finite expression in~\eqref{spc:eq:inverseexp} after increasing $A_M$. Choose $Y_M$ also so that $2\varepsilon_M(y)<1$. Then the two ceilings either agree or differ by one, giving at most two adjacent candidates.
\end{theorem}
\begin{proof}
Choose $B_M,N_0$ such that
$|\log z_n-F_M(n)|\le B_Mn^{-M-1}$ for $n\ge N_0$. On $[\xi/2,2\xi]$, for large $\xi$ we have $F_M'(x)\ge c\log\xi$. Choose $A_M$ large enough that this derivative times $\varepsilon_M$ dominates $B_Mn^{-M-1}$ throughout the interval. For every integer $n\ge\xi+\varepsilon_M$ there, $\log z_n\ge y$; for every integer $n<\xi-\varepsilon_M$ there, $\log z_n<y$. Integer nodes near $\xi/2$ and $2\xi$, together with monotonicity, extend these statements outside the interval once the finite initial segment is excluded by enlarging $Y_M$. Taking ceilings gives~\eqref{spc:eq:ceilings}. The inverse approximation error is of the same order as $\varepsilon_M$ and is absorbed by enlarging its constant. Finally $\varepsilon_M\to0$, so $Y_M$ can be enlarged until the interval width is less than $1$.
\end{proof}
The strict lower comparison and non-strict upper comparison match the definition with $\ge y$. Even an arbitrarily accurate asymptotic approximation need not have the correct ceiling at every real threshold near a jump. Thus no single-ceiling formula without an error envelope is asserted. The constants $A_M,Y_M$ have not been numerically evaluated or certified; the exact integer inverse in the code avoids them entirely within its declared resource bound.

\begin{remark}[{[write]} Section~\ref{spc:sec:inverse} in the terms of the transseries volume, 5~October 2026]\label{spc:rem:transseries}
The source cites the transseries collection~\cite{ProveItTransseries} for
``the Lambert core, residual transport, and staircase qualifications''.
Statement by statement, against the transseries volume
\path{Analysis/Transseries/docs/series-and-transseries/Transseries_And_Inversion/transseries_and_inversion.tex}:

\emph{(a) The core.} The core \eqref{spc:eq:core} is an instance of
\file{p0:prop:factorial-core}, not of \file{p0:thm:lambert-core}. The
equation $H(X)=y$ reads $X(\kappa'\log X+d')=y$ with $\kappa'=2$ and
$d'=-2\kappa$. The proposition gives $v=W_0\bigl((y/\kappa')e^{d'/\kappa'}\bigr)
=W_0\bigl(y/(2e^\kappa)\bigr)=w$, $X=\exp(v-d'/\kappa')=e^{\kappa+w}$ and core
slope $\kappa'(v+1)=2(1+w)=\Lambda$, exactly as in~\eqref{spc:eq:core}; its
side condition $\kappa'\log X+d'>0$ is $w>0$, which holds for $y>0$. The
volume's Lambert core proper concerns the linear--logarithmic phase
$aX+b\log X$; the phase $2X\log X$ here is outside its exponential--power
model, as the volume says of the double factorial.

\emph{(b) The recursion.} The recursion \eqref{spc:eq:inverseformal} is an
instance of \file{p0:thm:core-reversion} after the change of variable
$E=t\delta=\delta/X$ (the displacement measured in units of~$X$). Multiplying
\eqref{spc:eq:inverseformal} by $t$ and using
$\sum_{r\ge2}(-1)^ru^r/(r(r-1))=(1+u)\log(1+u)-u$ gives
\[
 \Lambda E+h(E)+f(t,E)=0,\qquad
 h(u)=2\bigl((1+u)\log(1+u)-u\bigr),
\]
\[
 f(t,u)=t\Bigl(\log K-\beta U-\beta\log(1+u)
 +\sum_{j\ge1}d_jt^j(1+u)^{-j}\Bigr),
\]
with $h\in u^2R[[u]]$ and $f\in tR[[t,u]]$ over the commutative
$\Q$-algebra $R=\Q[L,\log K,\log(2L),U][\Lambda^{-1}]$ of real rational
functions of~$U$, in which $\Lambda=2U-\log(2L)$ is invertible. The theorem's
unique solution $E=\sum_{n\ge1}e_nt^n$ has $e_{m+1}=v_m$, and its triangular
recurrence is the source's rule ``take minus the coefficient of $t^m$ \dots\
and divide by $\Lambda$''; for instance
$e_1=-\Lambda^{-1}[t^1]f=(\beta U-\log K)/\Lambda=v_0$. This is the same
shape as the volume's double-factorial instance
(\file{p0:rem:core-instances}), with the factor~$2$ in~$h$ and a different~$f$.

\emph{(c) The analytic step.} The analytic step of
Theorem~\ref{spc:thm:inverse} (a root of
$F_M=y$ near the finite expression, then division of the residual by the
slope) is the argument of \file{p0:thm:backward-error}, applied to $F_M$ on
$[X-1,X+1]$, where $F_M'\asymp\log X$; the source proves it directly.

\emph{(d) Sampled recovery.} The rounding in
Corollary~\ref{spc:cor:sampled} is \file{p0:thm:staircase}\,(3) for $A_n=\log z_n$ with $n_1=2$ (strictly
increasing and unbounded by~\eqref{spc:eq:increasing}), for every admissible
interpolation, since $\nu(\log z_n)=n$; the input bound
$|\xi_M(\log z_n)-n|<\tfrac12$ is the source's own estimate.

\emph{(e) The threshold enclosure.} Theorem~\ref{spc:thm:threshold} is
proved directly, by comparing
$\log z_n$ with $F_M$ at the integers; as proved it is not an instance of
\file{p0:thm:staircase}, which is stated for an interpolated inverse. Its
conclusion also follows from \file{p0:thm:staircase}\,(1) with one
interpolation chosen for the purpose (a second route, with possibly different
constants). \emph{Proof.} Put $r_n=\log z_n-F_M(n)$. By~\eqref{spc:eq:logzn}
there are $B_M$ and $N_0$ with $|r_n|\le B_Mn^{-M-1}$ for $n\ge N_0$. Since
$F_M'(x)\sim2\log x$, choose an integer $n_1\ge\max(N_0,3)$ with
$F_M'(x)\ge\log x$ for $x\ge n_1$ and $2^{M+1}B_Mn_1^{-M-1}<\log n_1$. On
$[n_1,\infty)$ let $\rho$ be the piecewise-linear function with $\rho(n)=r_n$
at the integers, and $\mathcal A=F_M+\rho$. Then $\mathcal A$ is continuous,
$\mathcal A(n)=\log z_n$, and on each $[n,n+1]$
\[
 \mathcal A'\ge\log n-|r_{n+1}-r_n|\ge\log n-2B_Mn^{-M-1}>0,
\]
because $2B_Mn^{-M-1}\le2^{M+1}B_Mn^{-M-1}$, which decreases in~$n$, while
$\log n$ increases. So $\mathcal A$ is strictly increasing and unbounded, an
admissible interpolation of \file{p0:def:three-inverses}, and for
$x\in[n,n+1]$, $|\mathcal A(x)-F_M(x)|\le B_Mn^{-M-1}\le2^{M+1}B_Mx^{-M-1}$.
For large $y$ put $\nu=\mathcal A^{-1}(y)$ and $\xi=\xi_M(y)$, both at least
$n_1$. Since $F_M(\xi)=y=\mathcal A(\nu)$, the mean value theorem gives
\[
 |\nu-\xi|\log\min(\nu,\xi)\le|F_M(\nu)-F_M(\xi)|
 =|F_M(\nu)-\mathcal A(\nu)|\le2^{M+1}B_M\nu^{-M-1}.
\]
The right side is at most $2^{M+1}B_Mn_1^{-M-1}<\log n_1\le\log\min(\nu,\xi)$,
so $|\nu-\xi|<1$; for $\xi\ge3$ this
gives $\nu^{-M-1}\le(\xi-1)^{-M-1}\le2^{M+1}\xi^{-M-1}$ and
$\log\min(\nu,\xi)\ge\log(\xi-1)\ge\tfrac12\log\xi$, hence
$|\nu-\xi|\le\varepsilon:=2^{2M+3}B_M\,\xi^{-M-1}/\log\xi$. For
$y\ge\log z_{n_1}$, \file{p0:thm:staircase}\,(1) gives
$\min\{n\ge n_1:\log z_n\ge y\}=\lceil\nu\rceil$, and this minimum is $N(y)$,
because $\log z_n<\log z_{n_1}\le y$ for $2\le n<n_1$
by~\eqref{spc:eq:increasing}. As the ceiling is nondecreasing,
$\lceil\xi-\varepsilon\rceil\le N(y)\le\lceil\xi+\varepsilon\rceil$,
which is~\eqref{spc:eq:ceilings} with $A_M=2^{2M+3}B_M$ for all sufficiently
large~$y$. \qed

The ordinary piecewise-linear interpolation of $\log z_n$ would not do
here: between the integers it differs from $F_M$ by about
$F_M''/8\approx1/(4x)$, which gives $|\nu-\xi|$ only of order
$1/(x\log x)$. When $2\varepsilon<1$, \file{p0:prop:certified-rounding}
(one residue class) says that at most two exact evaluations of $z_n$ decide
$N(y)$, which is the source's ``at most two adjacent candidates''. No
novelty is claimed for any of the inversion.
\end{remark}

\section{Exact computation and reproducibility}\label{spc:sec:computation}
The accompanying package has three deliberately separate layers.
\begin{enumerate}
\item The standard-library core computes exact integers, chain polynomials, fixed moments and cumulants, and a resource-bounded exact integer threshold inverse. It does not require a numerical value of $C$.
\item An optional SymPy program constructs kernel and correction coefficients by the finite recipe in Section~\ref{spc:sec:coefficients}, checks logarithmic and cumulant algebra, and checks formal inverse cancellation. Its output is exact symbolic algebra.
\item An optional mpmath program evaluates normalized sequences and centered moments at finite $n$. Its decimal outputs are diagnostics, not interval enclosures or proofs of error bounds.
\end{enumerate}
All public checks use explicit exceptions rather than assertions, so optimized Python execution does not remove the checks. Input lengths and resource parameters are validated before expensive operations. Computed large integers are converted to decimal in bounded chunks without changing Python's process-wide integer conversion limit. The package tests genuinely large output through the real command-line interface with \texttt{PYTHONINTMAXSTRDIGITS=640}. No global recursion, integer-conversion, or precision settings are changed; optional numerical work uses a local mpmath context.

\subsection{Independent finite checks}
Full partition-poset enumeration is feasible at small $n$ and provides a check independent of the rank-only recurrence. Enumerate all set partitions of $[n]$, test refinement for every pair, and sum chains over all strictly coarser partitions. This gives
\begin{align*}
 Z_1(q)&=1,\\
 Z_2(q)&=q,\\
 Z_3(q)&=q+3q^2,\\
 Z_4(q)&=q+13q^2+18q^3,\\
 Z_5(q)&=q+50q^2+205q^3+180q^4,\\
 Z_6(q)&=q+201q^2+1865q^3+4245q^4+2700q^5.
\end{align*}
The numbers of poset elements are $1,2,5,15,52,203$. These results also check the exact length convention.

For exact moment computations one may use ordinary derivatives $R_r(n)=Z_n^{(r)}(1)$:
\begin{equation}\label{spc:eq:derivrec}
 R_r(n)=\sum_{j<n}S(n,j)\bigl(R_r(j)+rR_{r-1}(j)\bigr),
\end{equation}
with $R_0(1)=1$, higher derivatives zero, and the second summand absent when $r=0$. Alternatively, Euler derivatives give raw moments directly by the binomial theorem applied to the additional first step. Stirling conversion relates these two implementations. Dividing the exact sums by $z_n$ gives rational moments, and the standard triangular moment--cumulant identity gives exact rational cumulants. The package checks both against the chain polynomials and the full poset enumeration.

A further exact check is the inequality $2^kk!S(n,n-k)\le(\fall nk)^2$ from Lemma~\ref{spc:lem:factorial}. Exact inverse checks verify both attainment and minimality: if the bounded search returns $n>2$, then $z_{n-1}<T\le z_n$ for its positive integer input $T$.

\subsection{Numerical diagnostics and their limits}
For reference, numerical substitution of $L=\log2$ in the exact coefficients gives
\[
 \begin{aligned}
 c_1&\approx 0.1155245300933242,& c_2&\approx 0.0109365636395843,\\
 c_3&\approx-0.0238254567258004,& c_4&\approx-0.0275946778045995.
 \end{aligned}
\]
Finite-$n$ corrected normalizations are consistent with
$C\approx1.098685805525187$. The following centered quantities from exact integer moment sums illustrate the smaller terms:
\[
 M_n=\E H_n-\frac n{2L}-\frac{\log n}{6}+\frac1{12n},\qquad
 V_n=\Var H_n-\frac{1-L}{4L^2}n+\frac{\log n}{12}-\frac1{24n}.
\]
\begin{center}
\begin{tabular}{rcc}
\toprule
$n$ & $M_n$ & $V_n$\\
\midrule
100 & $-0.56552859348282824923$ & $-0.22033553105629492173$\\
200 & $-0.56552929520692838409$ & $-0.22033372829125359404$\\
400 & $-0.56552946656943563565$ & $-0.22033328277962673835$\\
600 & $-0.56552949796082401647$ & $-0.22033320071484707943$\\
\bottomrule
\end{tabular}
\end{center}
These are compatible with $K_1\approx-0.565529523$ and $K_2\approx-0.220333135$. Neither these estimates nor the displayed estimate for $C$ is interval-certified. Their role is to detect algebraic or indexing mistakes; they are not used in any theorem proof or exact threshold decision.

The README gives the exact commands, supported bounds, dependency versions, source links, checksums, and deterministic rebuild procedure. The PDF has no creation date or trailer identifier; the archive uses a fixed member order, timestamp, permissions, and compression policy. A full rebuild is intended to preserve every archive member and the entire archive byte for byte under the recorded toolchain, in ordinary and optimized Python modes and with the supported integer conversion cap. Cross-version TeX byte identity is not promised. No third-party source PDFs or private research notes are included.

\begin{writenote}[The package here, and the intake reruns, 5~October 2026]
``The README'' and ``the archive'' above are the delivered ones. In this
report the four programs are \file{code/lengyel_exact.py},
\file{code/test_lengyel_exact.py}, \file{code/symbolic.py} and
\file{code/diagnostics.py} (delivered in \file{code/}); the rebuild tool
\file{code/build.py} was delivered at the package root; the recorded outputs
\file{generated/*} are \file{data/generated-*}; \file{SOURCES.json} and
\file{requirements-optional.txt} are in \file{data/}. The delivered PDF and
\file{SHA256SUMS} are not shipped, and the report README replaces the
delivered one. \file{code/build.py} expects the delivered layout
(\file{Report223.tex}, \file{SHA256SUMS}, \file{generated/}), so it does not
run in this directory; the report README gives a route from the arrival
commit's archive. The four programs need no layout: the test program finds
\file{lengyel_exact.py} beside itself. At placement, on copies, the three
programs regenerated \file{generated/*.txt} byte for byte, up to the CRLF line
ends of Windows standard output. At the write, on copies of the shipped
\file{code/} (Windows, Python 3.14.4; SymPy 1.14.0 and mpmath 1.3.0 as
pinned), the test program printed all eight PASS lines in 6~s, the symbolic
program its coefficients and three PASS lines in 106~s, and the diagnostics
program in 3~s; after CRLF to LF each output equals the shipped
\file{data/generated-*.txt} byte for byte. The full rebuild
(\file{build.py}), which compares pdfTeX builds and archive bytes under the
recorded TeX Live toolchain, was not run.
\end{writenote}

\section{Further questions within a clear boundary}\label{spc:sec:questions}
The arguments leave several concrete directions open without assuming their answers.
\begin{enumerate}
\item \textbf{Certified constants.} Turn the constructive bounds behind~\eqref{spc:eq:certificate} into practical interval enclosures for $C$, then extend them to derivatives of $\log C(e^t)$. This would support evaluated inverse enclosures rather than existence constants alone.
\item \textbf{Dependence on the order.} Quantify how $C_e$, the Taylor remainder constants, the finite patches, and the weighted contraction onset grow with $M$. Such estimates are prerequisites for a justified late-order or optimal-truncation analysis; the fixed-order theorem itself supplies none.
\item \textbf{A larger marking region.} The explicit disk $|q-1|\le1/16$ is convenient rather than optimal. A larger analytic domain would require tracking the absolute contraction and possible zeros of $C(q)$. The local proof does not establish a global marked-parameter result.
\item \textbf{Distributional refinements.} A local limit theorem, a quantitative normal approximation, or large deviations would require additional arguments not supplied here. Local limits and global large-deviation statements may also require control outside the small complex neighborhood. The cumulant expansion alone should not be promoted to those conclusions.
\item \textbf{Comparison with analytic iteration.} Relate the anchored recurrence limit and its marked derivatives directly to the contour constants in the prior analytic-iteration treatment, and clarify the full contents of the unrecovered Flajolet--Salvy manuscript if it becomes available.
\item \textbf{Aggregate versus depth asymptotics.} Compare this direct aggregate proof with consequences obtainable by summing proportional-depth iterated-Bell asymptotics, with particular attention to tail control and uniformity in depth. Their relationship deserves an actual derivation rather than an absence claim.
\end{enumerate}
These are research questions, not additional asserted theorems. The completed result of this report is the explicit fixed-order recurrence expansion, its local analytic marking, the cumulant consequences, and the specialized inverse bounds proved above.

\subsection{Further questions and research}\label{spc:sec:further}
\begin{writenote}[Added 5 October 2026, batch~106]
Following Vladimir Reshetnikov's standing rule of 4~October 2026, every
claim of the source that is not proved in full is stated here as an open
question, with its source, the argument offered and what is missing. The
six items above are Questions~\ref{spc:q:certified}--\ref{spc:q:depth}
below, in the same order; Questions~\ref{spc:q:inverse}
and~\ref{spc:q:variance} collect two statements the source makes elsewhere.
The intake found no false claim in the source; the one error it records is
in the OEIS entry (Remark~\ref{spc:rem:oeis}).
\end{writenote}
\begin{enumerate}
\item\label{spc:q:certified} \textbf{Certified constants} (source, item~1;
 Section~\ref{spc:sec:completion}: ``The proof is effective in principle''
 and ``The kernel remainder gives computable bounds, and finite-cutoff
 estimates make the weighted contraction effective''). Proved:
 \eqref{spc:eq:certificate} converts a verified bound
 $|d_n|\le A_Mn^{-M-2}$ into an enclosure of~$C$. Not supplied: any explicit
 $A_M$, onset or interval value of $C$, $K_1$, $K_2$ or the marked
 derivatives; the effectivity sentences are a sketch. The decimals
 $C\approx1.098685805525187$, $K_1\approx-0.565529523$,
 $K_2\approx-0.220333135$ are diagnostics (Section~\ref{spc:sec:computation};
 the dossier's recomputation agrees for~$C$ to about $5\cdot10^{-13}$). By
 Theorem~\ref{spc:thm:main}, $z_n/f_n=C\bigl(1+L/(6n)+O(n^{-2})\bigr)$, which
 answers the convergence-speed question at the end of Babai and Lengyel's
 Section~2~\cite{BabaiLengyel} (as the source reports it; not read by the
 write) asymptotically, but not with explicit constants. Missing: explicit
 kernel and contraction constants, then interval arithmetic.
\item\label{spc:q:order} \textbf{Dependence on the order} (source, item~2;
 Section~\ref{spc:sec:meaning}). Every statement is for fixed~$M$;
 convergence, Gevrey bounds, optimal truncation and exponentially small terms
 are not addressed. Missing: growth of $C_e$, the Cauchy constants, the
 patches and the contraction onset in~$M$.
\item\label{spc:q:disk} \textbf{A larger marking region} (source, item~3;
 Theorem~\ref{spc:thm:marked}). The disk $|q-1|\le1/16$ comes from the
 explicit contraction bound~\eqref{spc:eq:diskcontract}; the further shrink
 avoids zeros of $C(q)$, whose location is unknown. Missing: the contraction
 and the zero set of $C(q)$ on larger domains.

 \textup{[write, 9~October 2026, batch~138]} \emph{In part.} Part~IV of
 \file{a290354-iterated-euler-diagonals} (its Corollary~39.3,
 \file{ied:par:cor:marked}) proves, under the four inputs named in
 Question~\ref{spc:q:depth},
 $Z_n(q)\sim C(q)(n!)^2[2L(q)]^{-n}n^{-1-L(q)/3}$ for every real $q>0$,
 uniformly on compact positive $q$-intervals, with
 $C(q)=[2L(q)]^{L(q)/3}h(L(q))/(q+1)$ ($L(q)$ as here, $h$ as in that
 question). This $C(q)$ is positive on $q>0$, extends holomorphically to
 $\Re q>0$, and near $q=1$ is the $C(q)$ of Theorem~\ref{spc:thm:marked}, by
 the identity theorem. So the leading equivalent holds on the whole positive
 axis, and the analytic amplitude continues beyond the disk $|q-1|\le1/16$.
 Not given: the contraction or a fixed-order expansion beyond the disk, which
 that report asks for on neighbourhoods of positive compact intervals (its
 Section~48, Question~4, \file{ied:par:q:positive}); the zero set of $C(q)$
 on $\Re q>0$ (its Question~3, \file{ied:par:q:zeros}); and uniform complex
 asymptotics of $Z_n(q)$, which continuation of the amplitude alone does not
 give. Its Corollary~45.1 (\file{ied:par:cor:K}) expresses $K_1$ and $K_2$
 through $h(L)$, $h'(L)$ and $h''(L)$; no certified values follow (its
 Question~5).
\item\label{spc:q:distribution} \textbf{Distributional refinements} (source,
 item~4). A local limit theorem, a rate in the central limit theorem,
 Edgeworth terms or large deviations for $H_n$ are not proved; the cumulant
 expansion of Corollary~\ref{spc:cor:logmgf} is local in~$t$. Missing:
 control of $Z_n(q)$ away from $q=1$ (for a local limit theorem, on the whole
 unit circle).
\item\label{spc:q:analytic} \textbf{Comparison with analytic iteration, and
 priority} (source, item~5, and Section~\ref{spc:sec:prior}). As the
 source reports it, Prellberg's presentation~\cite{Prellberg} gives the
 leading equivalent with a contour expression for the constant; once the
 normalizations are matched, and provided that equivalent holds as stated,
 that expression and the recurrence limit~\eqref{spc:eq:C} are the same
 number, $\lim z_n/f_n$ (the write did not read the presentation). What is open is a direct derivation of one from
 the other, the marked derivatives
 $K_1,K_2$ in contour form, and the contents of the unfinished 1990
 Flajolet--Salvy manuscript (the OEIS attributes the decimal of $C$ to
 Flajolet and Salvy). Until then no first occurrence of $c_2,c_3,\dots$ is
 claimed. Missing: the sources.
\item\label{spc:q:depth} \textbf{Aggregate versus depth asymptotics}
 (source, item~6). The exact identity
 $z_n=\sum_{m\ge0}2^{-m-1}H(n,m)$ of
 \file{a139383-iterated-bell-diagonals} (Part~II, Section~17.2), where
 $H(n,m)=n![z^n](e^z-1)^{\circ m}$ are the iterated Bell numbers, links this
 report to the proportional-depth theorem \file{ibd:pd:thm:main} there. \emph{Heuristic} [write]: its leading form
 \file{ibd:pd:eq:factorial} reads
 $H(n,m)\approx(n-1)!\,2^{-n}m^n m^{-\beta/3}I(\beta)$ with $\beta=n/m$ and
 that report's amplitude $I$. The weight $2^{-m}m^n$ peaks at $m=n/L$
 ($\beta=L$, slope $\lambda=1/L\approx1.443$), with
 $\sum_m2^{-m}m^n\approx(n/(eL))^n\sqrt{2\pi n}/L$ by Laplace's method, and
 $n^ne^{-n}\sqrt{2\pi n}\sim n!$; formally, then,
 \[
  z_n\approx(n!)^2(2L)^{-n}n^{-1-L/3}\cdot\frac{L^{L/3-1}I(L)}2,
  \qquad\text{that is}\qquad
  C\overset{?}{=}\frac{L^{L/3-1}I(L)}2=\frac{D(1/L)}{2\sqrt{2\pi}\,L},
 \]
 with $D$ as in \file{ibd:pd:eq:amplitude}. The scale $f_n$, including the
 exponent $b=1+L/3$, comes out exactly. \emph{Observation, not a claim.} The
 write estimated $I(L)$ from exact values of $H(n,m)$ (by the recurrence)
 at the two integers $m$ nearest $n/L$, dividing by
 $(n-1)!\,2^{-n}m^nm^{-n/(3m)}$ and interpolating linearly in~$\beta$ to
 $\beta=L$. This gives $I(L)\approx1.66277$, $1.66098$, $1.66009$ and
 $L^{L/3-1}I(L)/2\approx1.102044$, $1.100861$, $1.100269$ at $n=200$, $300$,
 $400$; these exceed $C$ by $3.36$, $2.18$, $1.58$ thousandths, about
 $0.65/n$, and two-point extrapolation in $1/n$ gives $1.098494$ (from
 $n=200,300$) and $1.098492$ (from $n=300,400$), about $1.9\cdot10^{-4}$ below~$C$.
 The gap is the logarithmic correction of that theorem: by
 \file{ibd:pd:eq:P1} with $\delta=0$, $\beta=L$, the first correction
 $P_1/n$ contains $-\tfrac{L^2}{18}(\log n)^2/n$, which extrapolation in
 $1/n$ does not remove. Fitting $A+B/n+D_1\log n/n+D_2(\log n)^2/n$ to the
 sixteen values $n=100,120,\ldots,400$ (verification run, 5~October 2026)
 returns $D_2=-0.02937$, against the predicted $-L^2C/18=-0.02933$, and
 $A=1.098694$, within $10^{-5}$ of~$C$; with $D_2$ fixed at its predicted
 value, $A$ agrees with $C$ to $3\cdot10^{-6}$ ($n\ge200$). These are fits
 to numerics, not a proof. \textup{(Sentence replaced after the independent
 check of 5~October 2026; the first wording read ``about
 $1.9\cdot10^{-4}$ below~$C$, a gap attributable to the logarithmic
 corrections of that theorem''.)} Missing for a
 proof: that theorem is uniform only for slopes in compact sets and rests on
 the contour foundation of that report's Part~I, so the sum needs separate
 bounds for $m\le\varepsilon n$ and $m\ge n/\varepsilon$ (note
 $H(n,m)\le H(n,m+1)$, since $S(n,n)=1$); its corrections
 $P_j(\log n;\lambda,\delta)$ are polynomials in $\log n$ whereas
 \eqref{spc:eq:expansion} has none, which a summation would have to explain
 (at the level of $(\log n)^2/n$ it does: the tilt $m^{-n/(3m)}$ of the
 Laplace window contributes a factor
 $\exp\bigl(\tfrac{L^2}{18}(\log n)^2/n+\cdots\bigr)$, which cancels the
 term of $P_1$ above; the $\log n/n$ terms were not compared; added after
 the independent check of 5~October 2026);
 and $I(L)$ is not certified (that report certifies only $\beta=1$). If
 proved, the identity would express $C$ by that report's contour functional,
 which bears on Question~\ref{spc:q:analytic}. Conversely, nothing here gives
 a pointwise depth equivalent from the aggregate, as that report says.

 \textup{[write, 9~October 2026, batch~138]} \emph{The leading heuristic is
 now a theorem, under four unrefereed inputs.} Part~IV of
 \file{a290354-iterated-euler-diagonals} (its Theorems~39.1 and~39.2,
 Corollary~39.3 and equation~(39.12); labels \file{ied:par:thm:identity},
 \file{ied:par:thm:aggregate}, \file{ied:par:cor:marked},
 \file{ied:par:eq:Lengyel}) proves, uniformly for the rate $\ell$ of the
 geometric depth weight in compact subsets of $(0,\infty)$ (that report
 writes $L$ for it),
 \[
  \sum_{m\ge0}e^{-\ell m}H(n,m)\sim(n!)^2(2\ell)^{-n}n^{-1-\ell/3}
  \,\ell^{\ell/3-1}I(\ell).
 \]
 The depths outside the central window $|m-n/\ell|\le\sqrt n(\log n)^2$ are
 bounded by the rational composition majorant $H(n,m)\le n!(m/2)^{n-1}$ (its
 Lemma~43.1, from $(e^z-1)^{\circ m}\preceq z/(1-mz/2)$ coefficientwise).
 Through the exact mixture $z_n=\sum_m2^{-m-1}H(n,m)$ (its Lemma~44.1 at
 $q=1$) this gives
 \[
  C=\frac{L^{L/3-1}I(L)}2=\frac12(2L)^{L/3}h(L),\qquad L=\log2,
 \]
 the heuristic above, where $h$ is the density of that report's positive
 Laplace measure (not the $h(u)$ of Section~\ref{spc:sec:inverse}) and
 $I(\ell)=\ell\,2^{\ell/3}h(\ell)$ (its Theorem~39.1). The four inputs are
 Part~I of that report (the normalized Fatou coordinate
 \file{ied:lead:lem:fatou}, \file{ied:lead:eq:fatou} and the positive measure
 \file{ied:lead:prop:measure}) and, from
 \file{a139383-iterated-bell-diagonals}, the absolute transfer
 \file{ibd:pd:eq:absolute}, whose Part~II rests on its Part~I's interval
 certificate of common entry, and the holomorphy part of
 \file{ibd:pd:prop:positive}; all are unrefereed and unformalized, and no
 numerical value enters the proofs. Nothing proved in this report depends on
 them. Of what is missing above, the bounds outside compact slopes and the
 leading summation are now supplied; the summation of every correction, which
 would explain why the polynomials in $\log n$ of \file{ibd:pd:eq:P1} leave
 the logarithm-free $c_j(L)$ of~\eqref{spc:eq:expansion}, stays open (that
 report's Section~48, Question~2, \file{ied:par:q:summation}), and $I(L)$ is
 still not certified (its Question~5). With that report's exploratory density
 value $h(\log2)=2.03764226105\ldots$ (inverse Laplace, not interval
 certified) the right side is $1.0986858055251874\ldots$, within
 $4\cdot10^{-16}$ of the value $1.09868580552518701$ fitted from exact $z_n$
 (independent check, note after this list), and $I(L)=1.6577006128\ldots$.
 For the first correction, $\sum_m2^{-Lm}H(n,m)=2z_n$ exactly and the main
 term above is $2Cf_n$ at $\ell=L$, so under the same inputs
 Theorem~\ref{spc:thm:main} makes the ratio of the aggregate to its main term
 $1+L/(6n)+O(n^{-2})$ there, and Theorem~\ref{spc:thm:marked} gives
 $1+c_1(L(q))/n+O(n^{-2})$ at $\ell=L(q)$ for real $q$ near $1$. Beyond the
 disk, the write of Part~IV observed (uncertified) at $\ell=1$, that is
 $q=1/(e-1)$, $|q-1|>0.41$, that $n$ times the ratio minus one is $0.16669$ at
 $n=320$; this note's recomputation from the exact recurrence gives
 $0.1666941$, against $c_1(1)+c_2(1)/320=0.1666946$ from the polynomials
 of~\eqref{spc:eq:cs}. That is consistent with~\eqref{spc:eq:markedintro}
 holding with the same polynomials outside the disk, which is not proved.
\item\label{spc:q:inverse} \textbf{An evaluated inverse} (Section~\ref{spc:sec:inverse},
 after Theorem~\ref{spc:thm:threshold}: ``The constants $A_M,Y_M$ have not
 been numerically evaluated or certified''). Theorem~\ref{spc:thm:threshold}
 and Corollary~\ref{spc:cor:sampled} are existence statements; the shipped
 exact inverse is a bounded search. Missing: explicit $B_M$
 in~\eqref{spc:eq:logzn} and a certified $C$ (Question~\ref{spc:q:certified}).
\item\label{spc:q:variance} \textbf{Van Cutsem and Ycart's variance
 conjecture} (Section~\ref{spc:sec:prior}). Their general conjecture for a
 variance equivalent, under hypotheses broader than this model
 (\cite{VanCutsemYcart}, p.~997, as the source reports it), is not resolved;
 Corollary~\ref{spc:cor:logmgf} proves the variance expansion
 \eqref{spc:eq:varintro} for this model only. Missing: an argument for their
 general class.
\end{enumerate}
\begin{writenote}[Independent check of the write, 5~October 2026]
An adversarial check made by the intake after the write (commit
\file{44ebdfef9}) found every item valid and no mathematical error. It
confirmed the live entry A005121 (revision \#79) and its formula line
verbatim, and, paging all 79 revisions of the history, that the line first
appears in revision~\#8 (13~September 2003) as ``(Lengyel)'', that
revision~\#30 (18~April 2014) changed only the attribution to
``[Lengyel]'', and that the sign is unchanged since 2003. Chains counted
directly from the definition (refinement chains from the finest to the
coarsest set partition) give $z_1,\ldots,z_6=1,1,4,32,436,9012$; the
recurrence reproduces these and all 19 live data terms; and the printed line,
solved with generic coefficients, has the unique solution $-\mathcal Z$
through $z^7$. Applied to $\mathcal Z$ itself the printed right-hand side
equals $\mathcal Z-z$ exactly, so the line fails only at the coefficient
of~$z$ (wording of Remark~\ref{spc:rem:oeis} sharpened, with a dated note).
It confirmed the five parts of Remark~\ref{spc:rem:transseries} against the
transseries volume, including a line-by-line check of the second route to
Theorem~\ref{spc:thm:threshold} in part~(e). In
Question~\ref{spc:q:depth} every printed digit reproduces, and the
$1.9\cdot10^{-4}$ gap is explained by the term
$-\tfrac{L^2}{18}(\log n)^2/n$ of \file{ibd:pd:eq:P1}: a free fit recovers
$D_2=-0.02937$ against the predicted $-0.02933$, and the constant then
agrees with $C$ to about $10^{-5}$ (sentence replaced, with a dated note).
From exact $z_n$ to $n=400$, an eight-node fit
$z_n/f_n=C\bigl(1+\sum_{k\le7}a_kn^{-k}\bigr)$ on $n=225,\ldots,400$ gives
$C=1.09868580552518701$, and $a_1,\ldots,a_4$ agree with $c_1,\ldots,c_4$
of \eqref{spc:eq:cs} to $12$, $12$, $9$ and $7$ digits.
\end{writenote}

\begin{thebibliography}{99}
\bibitem{Lengyel}
T. Lengyel, \emph{On a recurrence involving Stirling numbers}, European Journal of Combinatorics \textbf{5} (1984), 313--321.
\href{https://doi.org/10.1016/S0195-6698(84)80035-9}{DOI 10.1016/S0195-6698(84)80035-9}.
\href{https://sites.oxy.edu/lengyel/originals/papers/hivatkozo/81989948\%20On\%20a\%20Recurrence\%20involving\%20Stirling\%20Numbers\%20EuropJComb1984.pdf}{Author-hosted full text}.

\bibitem{BabaiLengyel}
L. Babai and T. Lengyel, \emph{A convergence criterion for recurrent sequences with application to the partition lattice}, Analysis \textbf{12} (1992), 109--119.
\href{https://doi.org/10.1524/anly.1992.12.12.109}{DOI 10.1524/anly.1992.12.12.109}.
\href{https://sites.oxy.edu/lengyel/papers/analysis.pdf}{Author-hosted full text}.

\bibitem{Prellberg}
T. Prellberg, \emph{On the asymptotic analysis of a class of linear recurrences}, FPSAC presentation (2002).
\href{https://sites.oxy.edu/lengyel/papers/prellberg/recurrence.pdf}{Author-hosted mirror of the presentation}.
The final slide distinguishes formal higher corrections from the leading theorem.

\bibitem{Mishna}
M. Mishna, summary of T. Prellberg's 2002 seminar, \emph{On the asymptotic analysis of a class of linear recurrences}, Algorithms Seminar 2002--2004 (2005), 47--50.
\href{https://sites.oxy.edu/lengyel/papers/prellberg/Prellberg2002a.pdf}{Full text}.

\bibitem{VanCutsemYcart}
B. Van Cutsem and B. Ycart, \emph{Renewal-type behavior of absorption times in Markov chains}, Advances in Applied Probability \textbf{26} (1994), 988--1005.
\href{https://doi.org/10.2307/1427901}{DOI 10.2307/1427901}.
\href{https://sites.oxy.edu/lengyel/originals/papers/hivatkozo/Renewal-Type\%20Behavior\%20of\%20Absorption\%20Times\%20in\%20MCh\%20VanCutsem_Ycart\%20AdvApplProb94.pdf}{Author-archive full text}.

\bibitem{Finch}
S. R. Finch, \emph{Mathematical Constants}, Cambridge University Press (2003), Section 5.7, 316--321.
\href{https://sites.oxy.edu/lengyel/originals/MathConst2003.pdf}{Author-hosted scan}.
Printed pages 319--320 discuss the unfinished Flajolet--Salvy work and the chain-length limit law.

\bibitem{DickeyRosenberg}
E. H. Dickey and N. A. Rosenberg, \emph{Labeled histories and maximally probable labeled topologies with multifurcation}, arXiv:2511.16799; Discrete Applied Mathematics (2026).
\href{https://arxiv.org/abs/2511.16799}{arXiv record};
\href{https://doi.org/10.1016/j.dam.2026.04.032}{DOI 10.1016/j.dam.2026.04.032}.
The full arXiv version, rather than the published version, was the compared text.

\bibitem{OEIS}
OEIS Foundation, \emph{A005121}, Lengyel numbers.
\href{https://oeis.org/A005121}{Sequence entry};
\href{https://raw.githubusercontent.com/oeis/oeisdata/main/seq/A005/A005121.seq}{official raw source}.
The inspected source header is revision \#79, 30 May 2026, 16:39:47; its formal generating-function sign is discussed in Section~\ref{spc:sec:prior}.

\bibitem{ProveItBell}
V. Reshetnikov, ProveIt, \emph{Iterated Bell diagonals}, article in the OEIS sequence asymptotics collection, inspected at commit
\texttt{9f19e58de1fa83abdc8504c166975d724d8a9791}.
\href{https://github.com/VladimirReshetnikov/ProveIt/blob/9f19e58de1fa83abdc8504c166975d724d8a9791/SetTheory/Cardinals/docs/reports/generating-functions-and-asymptotics/oeis-sequence-asymptotics/a139383-iterated-bell-diagonals/article.tex}{Pinned article source}.
The compared blob is \texttt{ab5014ecd6366ef621d4d45393609bf177c88339}.
\par\textup{[write, 5 October 2026: that report has no author line (its title block reads ``A merged research report built from two manuscripts''). It is \file{a139383-iterated-bell-diagonals} in this collection, titled ``Asymptotics of iterated Bell diagonals (OEIS A139383, A261280)''; its article blob at the commit of this write is still \texttt{ab5014ec}.]}

\bibitem{ProveItTransseries}
V. Reshetnikov, ProveIt, \emph{Transseries and Inversion} and companion combinatorial material, repository revision
\texttt{9f19e58de1fa83abdc8504c166975d724d8a9791}.
\href{https://github.com/VladimirReshetnikov/ProveIt/tree/9f19e58de1fa83abdc8504c166975d724d8a9791/Analysis/Transseries/docs/series-and-transseries}{Pinned transseries collection}.
The Lambert core, residual transport, and staircase qualifications are used here as established general methods.
\end{thebibliography}
\end{document}
